\section{倫理原理・意思決定・組織倫理}

\subsection{平成24年度 問2}
\label{q:h24-reg-apt-02}

\noindent\textbf{原典：} sources/H24/aptitude.pdf、PDF物理1頁。

\textbf{公式正答：} 3（single）。

\paragraph{問題の要旨} 倫理綱領に沿う技術者の行動について、最も不適切な記述を選ぶ。

\Needspace{4\baselineskip}\paragraph{確認した条件・記述（独自要約）}\mbox{}\par

\noindent\textbf{方向：} 最も不適切な記述。\par

\noindent\textbf{対象：} 公衆、雇用者、顧客への責任を持つ技術専門職。\par

\paragraph{問われる核心} 公益確保と適切な情報開示

\paragraph{必要知識} 公衆への責任、秘密保持の正当理由、事実の確認、差別防止。

\paragraph{考え方・正答を選ぶ理由（査読待ち初稿）} 1・2・4・5は誠実性、公正、人格尊重、育成の行動である。3は公益を理由に公開の対象・時期・手続の検討を全て省いている。必要な警告や緊急対応は遅らせないが、危険の可能性があることと無条件で全面公開することは同じではない。3を選ぶ。

\Needspace{4\baselineskip}\paragraph{選択肢の検討（査読待ち初稿）}\mbox{}\par

\noindent\textbf{1：適切} 誠実な対応で専門職の品位・信頼を高める。。専門職の信用は、能力とともに誠実な履行によって支えられる。\par

\noindent\textbf{2：適切} 不当な対価を直接にも間接にも与えず、求めず、受け取らない。。不当な利益の授受は判断の公正さを損なう。直接・間接を問わないという趣旨は妥当である。\par

\noindent\textbf{3：不適切} 公衆の安全・健康・福祉又は環境へ害の可能性があれば、即時・無条件に情報を公開する。。危険の確認、被害の緊急性、適切な通報先、秘密・個人情報の扱いを検討せず、あらゆる情報を無条件で公表することは適切でない。\par

\noindent\textbf{4：適切} 人種、性、年齢、地位、所属、思想・宗教等によって差別せず、人権と人格を尊重する。。人格尊重と公平な取扱いは倫理的な協業の基礎である。\par

\noindent\textbf{5：適切} 専門知識と経験を生かして、次世代の技術者・研究者を育てる。。専門職の知識継承と人材育成に沿う。\par

\paragraph{再利用する考え方} 事実と価値判断、行為の結果と手続を区別し、公益確保と適切な情報開示から判断の理由を説明する。

\paragraph{共通理論} \hyperref[theory:ethics:public]{第\ref*{theory:ethics:public}節}, \hyperref[theory:ethics:disclosure]{第\ref*{theory:ethics:disclosure}節}

\par 参考：\cite{ipej_ethics2011,ipej_ethics2023}。

\paragraph{関連問題} \hyperref[q:r02-reg-apt-02]{R02-REG-APT-02}

\par 難易度 2（校準初稿）、優先度未評価。

\subsection{平成25年度 問12}
\label{q:h25-reg-apt-12}

\noindent\textbf{原典：} sources/H25/aptitude.pdf、PDF物理13頁。

\textbf{公式正答：} 5（single）。

\paragraph{問題の要旨} 予算制約下のプロジェクト運営で、透明性を備えた最も適切な対応を選ぶ。

\Needspace{4\baselineskip}\paragraph{確認した条件・記述（独自要約）}\mbox{}\par

\noindent\textbf{条件：} 経済性とともに活動の透明性に注意する。\par

\noindent\textbf{E：} 実験の遅れで予定期間内の完了が難しくなり、上司へ報告し了承された。\par

\paragraph{問われる核心} 予算と進捗の説明責任

\paragraph{必要知識} 予算承認、費用と効果、重要性の定性的評価、記録と報告。

\paragraph{考え方・正答を選ぶ理由（査読待ち初稿）} 予算は「使い切る」ことを目的とせず、目的に必要な支出を承認された範囲で行う。5は進捗の悪化を共有し、残額を記録して返上する。1と4は手続を省略し、2と3は評価軸を恣意的に狭めるため、最も適切な5を選ぶ。

\Needspace{4\baselineskip}\paragraph{選択肢の検討（査読待ち初稿）}\mbox{}\par

\noindent\textbf{1：設問条件に不一致} 良い成果が出るとの確信から予算超過を上司へ報告せず遂行する。。結果への自信は、予算変更の承認と報告を省略する根拠にならない。\par

\noindent\textbf{2：設問条件に不一致} 事業継続計画の複数の危機に、対策が安価なものから順に対応する。。安価さだけでは危機の重大性、発生可能性、重要業務への影響と対策の効果を評価できない。\par

\noindent\textbf{3：設問条件に不一致} 費用対効果の順位で計数できない項目は除外し、それを考慮する上司の指示を拒否する。。金銭へ換算しにくい安全、信頼、社会的影響も意思決定に必要である。評価法を明示して扱う。\par

\noindent\textbf{4：設問条件に不一致} 継続のため必要機器を私費で購入し、予算超過への問題はないとする。。資金の出所を私費にしても、支出・所有・管理を含む透明な組織手続が不要になるわけではない。\par

\noindent\textbf{5：設問条件に一致} 遅延を報告して了承を得たうえで、余る予算を未消化分として計上し返上する。。遅延の共有と了承、未消化額の記録、不要な支出の回避が透明性に沿う。\par

\paragraph{再利用する考え方} 事実と価値判断、行為の結果と手続を区別し、予算と進捗の説明責任から判断の理由を説明する。

\paragraph{共通理論} \hyperref[theory:ethics:procedure]{第\ref*{theory:ethics:procedure}節}

\par 参考：\cite{ipej_ethics2023}。

\par 難易度 2（校準初稿）、優先度未評価。

\subsection{平成27年度 問6}
\label{q:h27-reg-apt-06}

\noindent\textbf{原典：} sources/H27/aptitude.pdf、PDF物理6頁。

\textbf{公式正答：} 3（single）。

\paragraph{問題の要旨} 功利主義・義務・徳などの倫理原理について、最も不適切な記述を選ぶ。

\Needspace{4\baselineskip}\paragraph{確認した条件・記述（独自要約）}\mbox{}\par

\noindent\textbf{前提：} 複数の倫理原理が並び立つ。人を単に技術のための手段と扱わないという倫理判断から導入する。\par

\paragraph{問われる核心} 結果・義務・徳による倫理判断

\paragraph{必要知識} 帰結主義と功利主義、義務と人格尊重、徳、負担の分配。

\paragraph{考え方・正答を選ぶ理由（査読待ち初稿）} 複数の原理は違う問いを立てる。全体への影響、侵害してはならない権利、専門職の人格をそれぞれ調べることで、直観を点検できる。技術者倫理の共通見解として唯一の原理を要求する3が誤りである。5は功利主義への批判を述べたもので、少数者の犠牲を本教材が推奨する意味ではない。

\Needspace{4\baselineskip}\paragraph{選択肢の検討（査読待ち初稿）}\mbox{}\par

\noindent\textbf{1：適切} 直観的判断の根拠へ戻り、倫理原理で批判的に再検討できる。。直観を結論の根拠として固定せず、理由を検討する役割は適切である。\par

\noindent\textbf{2：適切} 最大多数の最大幸福、厳守すべき義務、義務以上の有徳さへ注目するなどで原理に違いがある。。結果、義務、人格という評価対象の違いを述べている。\par

\noindent\textbf{3：不適切} 諸原理は補完し合うが、唯一正しい原理が存在すべきとの共通見解が技術者倫理にある。。各原理の観点が異なることから、単一の正しい原理を専門職全体の合意として断定できない。\par

\noindent\textbf{4：適切} 功利主義は動機・有徳さより行為の結果と幸福を重視する。。功利主義は結果を重視する帰結主義の一形態である。\par

\noindent\textbf{5：適切} 功利主義には全体の利益のため少数の不利益を我慢させやすいという指摘がある。。集計上の総利益と、少数者の権利・負担の公平性が衝突し得るという批判である。\par

\paragraph{再利用する考え方} 事実と価値判断、行為の結果と手続を区別し、結果・義務・徳による倫理判断から判断の理由を説明する。

\paragraph{共通理論} \hyperref[theory:ethics:principles]{第\ref*{theory:ethics:principles}節}

\par 参考：\cite{scu_ethics_framework}。

\par 難易度 2（校準初稿）、優先度未評価。

\subsection{平成29年度 問13}
\label{q:h29-reg-apt-13}

\noindent\textbf{原典：} sources/H29/aptitude.pdf、PDF物理12頁。

\textbf{公式正答：} 3（single）。

\paragraph{問題の要旨} 黄金律による他者への配慮と技術者の優先責任を表す空欄の組を選ぶ。

\Needspace{4\baselineskip}\paragraph{確認した条件・記述（独自要約）}\mbox{}\par

\noindent\textbf{A：} 自分が他者から与えられる何かを立場交換で考える。\par

\noindent\textbf{B：} 自分が受けることをどう判断するか。\par

\noindent\textbf{C・D：} 公衆の何を優先するか。\par

\noindent\textbf{並び：} A／B／C／D。\par

\paragraph{問われる核心} 黄金律の可逆性と公衆への優先責任

\paragraph{必要知識} 立場交換、他者の価値観、安全・健康、倫理綱領。

\paragraph{考え方・正答を選ぶ理由（査読待ち初稿）} 黄金律に基づく可逆性の点検は、自分が不利益を受ける側ならその扱いを許容できるかという問いである。技術者の公衆への優先責任は安全・健康で、3となる。自分だけが許容したことを全他者の同意へ置き換えず、異なる事情・価値観も聞く。人格や環境が重要でないという結論も導かない。

\Needspace{4\baselineskip}\paragraph{選択肢の検討（査読待ち初稿）}\mbox{}\par

\noindent\textbf{1：設問条件に不一致} 不利益／無視／安全／人格。立場交換で不利益の許容を問うのでBが違い、優先責任のDも合わない。\par

\noindent\textbf{2：設問条件に不一致} 不道徳／許容／環境／人格。Aは不利益、C・Dは安全・健康の組である。\par

\noindent\textbf{3：設問条件に一致} 不利益／許容／安全／健康。受け手側で不利益を許容できるか考え、公衆の安全・健康を優先する。\par

\noindent\textbf{4：設問条件に不一致} 不道徳／無視／環境／健康。A・B・Cが文脈と合わない。\par

\noindent\textbf{5：設問条件に不一致} 不利益／許容／環境／人格。A・Bは合うがC・Dは安全・健康である。\par

\paragraph{再利用する考え方} 黄金律の可逆性と公衆への優先責任について、対象・要件・例外・時点を分け、類似用語や無条件の断定を避ける。

\paragraph{共通理論} \hyperref[theory:ethics:tests]{第\ref*{theory:ethics:tests}節}, \hyperref[theory:ethics:public]{第\ref*{theory:ethics:public}節}

\par 参考：\cite{davis_seven_steps,ipej_ethics2011}。

\par 難易度 3（校準初稿）、優先度未評価。

\subsection{平成29年度 問15}
\label{q:h29-reg-apt-15}

\noindent\textbf{原典：} sources/H29/aptitude.pdf、PDF物理14頁。

\textbf{公式正答：} 2（single）。

\paragraph{問題の要旨} 倫理的意思決定の促進要因と阻害要因を対比した7行のうち、不適切な対比の数を選ぶ。

\Needspace{4\baselineskip}\paragraph{確認した条件・記述（独自要約）}\mbox{}\par

\noindent\textbf{ア：} 促進=利他主義、阻害=利己主義。\par

\noindent\textbf{イ：} 促進=希望・勇気、阻害=失望・おそれ。\par

\noindent\textbf{ウ：} 促進=正直・誠実、阻害=自己ぎまん。\par

\noindent\textbf{エ：} 促進=知識・専門能力、阻害=無知。\par

\noindent\textbf{オ：} 促進=公共的志向、阻害=自己中心的志向。\par

\noindent\textbf{カ：} 促進=指示・命令への批判精神、阻害=無批判な受入れ。\par

\noindent\textbf{キ：} 促進=依存的思考、阻害=自律的思考。\par

\paragraph{問われる核心} 道徳的自律と無批判な依存

\paragraph{必要知識} 利他性、自己欺瞞、知識不足、指示の吟味、相談と責任放棄の違い。

\paragraph{考え方・正答を選ぶ理由（査読待ち初稿）} ア〜カは他者と公衆への配慮、勇気、事実への誠実さ、能力、指示の吟味を促進側へ置く。キだけは自分で理由を検討する自律を阻害側へ置き、他者への依存を促進側へ置いて逆になっている。不適切1件で2を選ぶ。協力や専門家への相談を「依存的思考」と同一視しない。

\Needspace{4\baselineskip}\paragraph{選択肢の検討（査読待ち初稿）}\mbox{}\par

\noindent\textbf{1：設問条件に不一致} 0。キで促進と阻害が逆なので0ではない。\par

\noindent\textbf{2：設問条件に一致} 1。キだけが逆の対比で、不適切なものは1件。\par

\noindent\textbf{3：設問条件に不一致} 2。ア〜カに追加の不適切行を数える根拠がない。\par

\noindent\textbf{4：設問条件に不一致} 3。逆の行はキだけで3件ではない。\par

\noindent\textbf{5：設問条件に不一致} 4。逆の行はキだけで4件ではない。\par

\paragraph{再利用する考え方} 事実と価値判断、行為の結果と手続を区別し、道徳的自律と無批判な依存から判断の理由を説明する。

\paragraph{共通理論} \hyperref[theory:ethics:obstacles]{第\ref*{theory:ethics:obstacles}節}

\par 参考：\cite{davis_seven_steps}。

\paragraph{関連問題} \hyperref[q:r07-reg-apt-03]{R07-REG-APT-03}

\par 難易度 2（校準初稿）、優先度未評価。

\subsection{平成29年度 問3}
\label{q:h29-reg-apt-03}

\noindent\textbf{原典：} sources/H29/aptitude.pdf、PDF物理3頁。

\textbf{公式正答：} 3（single）。

\paragraph{問題の要旨} 納期を短縮し原価を下げる工程変更の提案を検討する際、4事項の優先順位を選ぶ。

\Needspace{4\baselineskip}\paragraph{確認した条件・記述（独自要約）}\mbox{}\par

\noindent\textbf{状況：} 材料の納入不足で遅延が生じ、納入業者が工程の一部変更を条件に早く安く納入できると提案。\par

\noindent\textbf{表の順：} 1番→2番→3番→4番。\par

\paragraph{問われる核心} 安全と品質を制約にした工程変更

\paragraph{必要知識} 公衆の安全優先、要求仕様、変更管理、費用と納期の比較。

\paragraph{考え方・正答を選ぶ理由（査読待ち初稿）} 工程を変えても安全と要求品質が確保されるかを先に確認する。その制約を満たす案について納期と原価を比較するので3である。この設問の順位を、全ての契約で納期と費用の関係が常に固定されるという一般法則へ拡張しない。

\Needspace{4\baselineskip}\paragraph{選択肢の検討（査読待ち初稿）}\mbox{}\par

\noindent\textbf{1：設問条件に不一致} 納期→原価→品質→安全。納期と費用を守るため安全と品質の確認を後回しにしている。\par

\noindent\textbf{2：設問条件に不一致} 安全→原価→品質→納期。安全を先にする点はよいが、仕様・性能へ関わる品質より原価を先にしている。\par

\noindent\textbf{3：設問条件に一致} 安全→品質→納期→原価。この提案の条件ではまず安全と必要品質を確保し、残る実行可能な案の納期と費用を検討する。\par

\noindent\textbf{4：設問条件に不一致} 品質→納期→安全→原価。安全を納期より下にしている。\par

\noindent\textbf{5：設問条件に不一致} 品質→安全→原価→納期。品質を安全より先にし、残る二事項もこの候補比較での適切な順序と異なる。\par

\paragraph{再利用する考え方} 事実と価値判断、行為の結果と手続を区別し、安全と品質を制約にした工程変更から判断の理由を説明する。

\paragraph{共通理論} \hyperref[theory:ethics:public]{第\ref*{theory:ethics:public}節}, \hyperref[theory:ethics:procedure]{第\ref*{theory:ethics:procedure}節}

\par 参考：\cite{ipej_ethics2011}。

\par 難易度 2（校準初稿）、優先度未評価。

\subsection{平成29年度 問7}
\label{q:h29-reg-apt-07}

\noindent\textbf{原典：} sources/H29/aptitude.pdf、PDF物理7頁。

\textbf{公式正答：} 4（single）。

\paragraph{問題の要旨} 新技術・新工法の品質問題に対する関係者の対応で、最も適切なものを選ぶ。

\Needspace{4\baselineskip}\paragraph{確認した条件・記述（独自要約）}\mbox{}\par

\noindent\textbf{状況：} 公共施設等でデータ虚偽報告が発生。開発者、現場、管理者、評価機関、発注者に専門能力と共通認識が必要。\par

\paragraph{問われる核心} 個人の誠実性と組織の再発防止

\paragraph{必要知識} 原記録と報告の真正性、ガバナンス、専門能力の共有、比例した対応。

\paragraph{考え方・正答を選ぶ理由（査読待ち初稿）} データの真正性を守り、現場と管理者が問題を共有し、評価機関も含めて専門能力を改善する必要がある。4はその学習と再発防止を進める。1・2は責任の転嫁、3は手続の検討不足、5は一律の過大対応であり、4を選ぶ。

\Needspace{4\baselineskip}\paragraph{選択肢の検討（査読待ち初稿）}\mbox{}\par

\noindent\textbf{1：設問条件に不一致} 現場で異常が反復するが、開発担当者の電話指示どおり日常的にデータを修正して提出する。。担当者の権威は、異常の原因を調べず原記録と異なる報告をする根拠にならない。\par

\noindent\textbf{2：設問条件に不一致} 支店で対策会議を開くが期待を伝え、現場だけで解決してから報告するよう指示する。。会議開催だけでは十分でなく、必要な組織的専門能力と責任を現場へ押し付けている。\par

\noindent\textbf{3：設問条件に不一致} 内部対策フローはあるが組織の倫理欠如を懸念し、会議前日にマスコミへ告発する。。緊急性、内部手続の実効性、告発先等を検討した記述がなく、この条件では会議前の報道告発を最適とできない。告発を一律に否定する趣旨ではない。\par

\noindent\textbf{4：設問条件に一致} 評価機関・団体が横断連携し、類似技術の問題と知見の共有・研修を進める。。個別現場の経験を組織・業界の知識として共有し、共通の専門能力を高める再発防止に沿う。\par

\noindent\textbf{5：設問条件に不一致} 民間開発の新技術を全て中止し、全従来工法で巧妙な偽装を前提とする抜き打ち立会検査を標準導入する。。検査強化自体は有用となり得るが、全技術を一律中止する措置は個別のリスク評価を欠き、過大な禁止となる。\par

\paragraph{再利用する考え方} 事実と価値判断、行為の結果と手続を区別し、個人の誠実性と組織の再発防止から判断の理由を説明する。

\paragraph{共通理論} \hyperref[theory:ethics:organization]{第\ref*{theory:ethics:organization}節}, \hyperref[theory:ethics:disclosure]{第\ref*{theory:ethics:disclosure}節}

\par 参考：\cite{ipej_ethics2023,scj_code2013}。

\par 難易度 2（校準初稿）、優先度未評価。

\subsection{平成30年度 問15}
\label{q:h30-reg-apt-15}

\noindent\textbf{原典：} sources/H30/aptitude.pdf、PDF物理15頁。

\textbf{公式正答：} 4（single）。

\paragraph{問題の要旨} 公務員倫理の範囲と、予防倫理・志向倫理を区別し、5空欄の組合せを選ぶ。

\Needspace{4\baselineskip}\paragraph{確認した条件・記述（独自要約）}\mbox{}\par

\noindent\textbf{ア：} 許認可・補助金等の相手方で、贈与・接待・割り勘のゴルフや旅行などの規制対象。\par

\noindent\textbf{イ：} 倫理規程で示す公務員倫理の狭い範囲。\par

\noindent\textbf{ウ：} 望ましい行動や姿勢まで含む公務員倫理の範囲。\par

\noindent\textbf{エ：} するなという服務規律を典型とする消極的倫理。\par

\noindent\textbf{オ：} した方がよいことへ取り組む積極的倫理。\par

\noindent\textbf{時点：} 平成30年度の制度と文章の語句を確認。現在の具体的な適用・例外は人事院資料を別途確認する。\par

\paragraph{問われる核心} 予防倫理と志向倫理の補完

\paragraph{必要知識} 利害関係者、倫理法令の対象、狭義/広義、予防と積極的な社会貢献。

\paragraph{考え方・正答を選ぶ理由（査読待ち初稿）} 職務との利害関係を持つ相手を利害関係者と呼ぶ。法・規程による限定された規制は狭義、望ましい姿勢も含めると広義である。悪い行為を防ぐ側面は予防、よりよい行為へ向かう側面は志向で、4を選ぶ。規程の詳細や例外をこの要約だけで実務へ適用しない。

\Needspace{4\baselineskip}\paragraph{選択肢の検討（査読待ち初稿）}\mbox{}\par

\noindent\textbf{1：設問条件に不一致} 利害関係者／狭義／広義／規律／自律。対象と範囲は合うが、対になる技術者倫理の名称は予防と志向である。\par

\noindent\textbf{2：設問条件に不一致} 知人／狭義／広義／予防／自律。アは単なる知人ではなく利害関係者であり、オも志向が適切。\par

\noindent\textbf{3：設問条件に不一致} 知人／広義／狭義／規律／志向。規制対象と広狭の関係が違い、エの名称も違う。\par

\noindent\textbf{4：設問条件に一致} 利害関係者／狭義／広義／予防／志向。対象、範囲、消極/積極の5対応が全て合う。\par

\noindent\textbf{5：設問条件に不一致} 利害関係者／広義／狭義／予防／自律。広狭が逆で、積極倫理の名称が志向になっていない。\par

\paragraph{再利用する考え方} 事実と価値判断、行為の結果と手続を区別し、予防倫理と志向倫理の補完から判断の理由を説明する。

\paragraph{共通理論} \hyperref[theory:ethics:aspiration]{第\ref*{theory:ethics:aspiration}節}

\par 参考：\cite{jinji_ethics2017,jinji_ethics_qa,fudano_ethics2021}。

\par 難易度 2（校準初稿）、優先度未評価。

\subsection{平成30年度 問4}
\label{q:h30-reg-apt-04}

\noindent\textbf{原典：} sources/H30/aptitude.pdf、PDF物理4頁。

\textbf{公式正答：} 1（single）。

\paragraph{問題の要旨} 倫理規程と情報開示、継続学習、法規制外の危険対応、組織文化、倫理的能力を判定する。

\Needspace{4\baselineskip}\paragraph{確認した条件・記述（独自要約）}\mbox{}\par

\noindent\textbf{方向：} 最も不適切。\par

\noindent\textbf{時点：} 平成30年。背景の2011年倫理綱領と現在の2023年指針を区別。\par

\paragraph{問われる核心} 倫理規程を行動へ適用するための条件と能力

\paragraph{必要知識} 倫理規程を行動へ適用するための条件と能力。倫理規程を守ることと、全ての情報を条件なしに公表することは同一ではない。1の即時・無条件という限定が不適切である。緊急の危害防止を遅らせず、必要な警告と適切な通報を具体的に選ぶ。2〜5は継続学習・公益・対話・倫理能力という責任を支える。

\paragraph{考え方・正答を選ぶ理由（査読待ち初稿）} 倫理規程を守ることと、全ての情報を条件なしに公表することは同一ではない。1の即時・無条件という限定が不適切である。緊急の危害防止を遅らせず、必要な警告と適切な通報を具体的に選ぶ。2〜5は継続学習・公益・対話・倫理能力という責任を支える。

\Needspace{4\baselineskip}\paragraph{選択肢の検討（査読待ち初稿）}\mbox{}\par

\noindent\textbf{1：不適切} 倫理規程に抵触する可能性があるときは即時・無条件に公開する。。危険の緊急性、事実確認、通報先、必要な範囲、無関係な秘密や個人情報を考慮せず、可能性だけで無条件公開する一般則は不適切。\par

\noindent\textbf{2：適切} 学協会等の組織外の講習にも参加し、能力・資質の向上を図る。。知識の陳腐化を防ぎ、組織内では得られない専門的な検討にも接する継続学習の趣旨である。\par

\noindent\textbf{3：適切} 法規制がなくても公衆への危険を察知したら対応する責務がある。。法規制は全危険を事前に網羅しない。専門家として必要な警告や予防措置を検討する。\par

\noindent\textbf{4：適切} 倫理上の問題を自由に議論し行動できる組織文化を育てる。。異論や相談を封じる組織では不正や危険が見落とされやすい。規程と実際の行動環境を結び付ける。\par

\noindent\textbf{5：適切} 専門学識に加え倫理的行動のための能力も資質能力に含む。。知識を持つだけでは責任ある使用を保証しない。事実と価値の判断、他者への説明・行動の能力が必要。\par

\paragraph{再利用する考え方} 対象、条件、時点を先に確認する。倫理規程を行動へ適用するための条件と能力を具体的な事例へ適用し、技術的原因・法的要件・倫理的判断のどれを問うか区別する。

\paragraph{共通理論} \hyperref[theory:ethics:disclosure]{第\ref*{theory:ethics:disclosure}節}, \hyperref[theory:ethics:organization]{第\ref*{theory:ethics:organization}節}, \hyperref[theory:ethics:public]{第\ref*{theory:ethics:public}節}

\par 参考：\cite{ipej_ethics2011,ipej_ethics2023,mext_competency2014}。

\par 難易度 2（校準初稿）、優先度未評価。

\subsection{平成30年度 問5}
\label{q:h30-reg-apt-05}

\noindent\textbf{原典：} sources/H30/aptitude.pdf、PDF物理5頁。

\textbf{公式正答：} 4（single）。

\paragraph{問題の要旨} 工学系学会の行動規範の前文に入る5語の組を選ぶ。

\Needspace{4\baselineskip}\paragraph{確認した条件・記述（独自要約）}\mbox{}\par

\noindent\textbf{A：} 未来につながる社会の性質。\par

\noindent\textbf{B：} 未来世代の権利。\par

\noindent\textbf{C・D：} 人類・社会に貢献すべき対象。\par

\noindent\textbf{E：} 専門家が協力する相手の見識。\par

\noindent\textbf{並び：} A／B／C／D／E。電気学会2007年行動規範の前文に対応。\par

\paragraph{問われる核心} 持続可能性・世代間責任・専門能力の限界

\paragraph{必要知識} 学会行動規範、未来世代、安全と健康、多様な見識、専門職の協同。

\paragraph{考え方・正答を選ぶ理由（査読待ち初稿）} 持続可能な社会を構築し、未来世代の生存権を確保し、安全と健康に貢献し、多様な見識を結集するという前文の構造から4である。電気技術の成果だけでなく、国際協調・権利・専門能力の限界を認識する文脈で読む。2021年改訂の倫理綱領と2007年行動規範を同一年の資料とはしない。

\Needspace{4\baselineskip}\paragraph{選択肢の検討（査読待ち初稿）}\mbox{}\par

\noindent\textbf{1：設問条件に不一致} 持続／生存／安全／健康／同様。Eは同じ見識の人だけではなく、多様な見識の人々である。\par

\noindent\textbf{2：設問条件に不一致} 持続／幸福／安定／安心／同様。B・C・D・Eが前文の保護対象・協力の意味と違う。\par

\noindent\textbf{3：設問条件に不一致} 進歩／幸福／安定／安心／同様。5語とも対応しない。技術進歩だけでは持続可能性を保証できない。\par

\noindent\textbf{4：設問条件に一致} 持続／生存／安全／健康／多様。持続可能性、未来世代の生存、安全・健康、多様な見識との協同が対応する。\par

\noindent\textbf{5：設問条件に不一致} 進歩／幸福／安全／安心／多様。C・Eは合うが、A・B・Dが違う。\par

\paragraph{再利用する考え方} 持続可能性・世代間責任・専門能力の限界について、対象・要件・例外・時点を分け、類似用語や無条件の断定を避ける。

\paragraph{共通理論} \hyperref[theory:ethics:code]{第\ref*{theory:ethics:code}節}

\par 参考：\cite{ieej_conduct2007}。

\par 難易度 3（校準初稿）、優先度未評価。

\subsection{令和元年度再試験 問3}
\label{q:r01-ret-apt-03}

\noindent\textbf{原典：} sources/R01\_retest/aptitude.pdf、PDF物理3頁。

\textbf{公式正答：} 1（single）。

\paragraph{問題の要旨} 倫理規程・行動規範を制定する狙いに、5記述が含まれるかを判定する。

\Needspace{4\baselineskip}\paragraph{確認した条件・記述（独自要約）}\mbox{}\par

\noindent\textbf{ア：} 一般社会と組織との契約について明確に意思表示する。\par

\noindent\textbf{イ：} 組織のメンバーの目指す理想を表す。\par

\noindent\textbf{ウ：} 倫理的行動の実践的な指針を示す。\par

\noindent\textbf{エ：} 将来のメンバーの教育に使う。\par

\noindent\textbf{オ：} 組織自体の在り方を議論する機会を与える。\par

\noindent\textbf{並び：} ア・イ・ウ・エ・オの順で○/×。\par

\paragraph{問われる核心} 倫理規程の社会的・教育的機能

\paragraph{必要知識} 倫理プログラム、理想と実践、教育、組織の継続的な自己点検。

\paragraph{考え方・正答を選ぶ理由（査読待ち初稿）} 倫理規程は違反者を罰するだけの文書ではない。社会と組織の信頼関係を示し、行動指針と教育の基準を共有し、制定・見直しの議論を通じて組織の目的を点検する。ア〜オは全て狙いに含まれ、1を選ぶ。「契約」はここでは社会的な約束の意味であり、個別の法的契約の成立を断定しない。

\Needspace{4\baselineskip}\paragraph{選択肢の検討（査読待ち初稿）}\mbox{}\par

\noindent\textbf{1：設問条件に一致} ○○○○○。5項目とも社会への説明、理想、実践、教育、組織の自己点検という目的に含まれる。\par

\noindent\textbf{2：設問条件に不一致} ○○○○×。制定の議論は組織の在り方を点検する機会でもあり、オを除外できない。\par

\noindent\textbf{3：設問条件に不一致} ○×○○○。望ましい理想を表明するイを除外できない。\par

\noindent\textbf{4：設問条件に不一致} ○○×○○。行動の実践的指針を示すウを除外できない。\par

\noindent\textbf{5：設問条件に不一致} ○○○×○。新たなメンバーを教育するエを除外できない。\par

\paragraph{再利用する考え方} 事実と価値判断、行為の結果と手続を区別し、倫理規程の社会的・教育的機能から判断の理由を説明する。

\paragraph{共通理論} \hyperref[theory:ethics:organization]{第\ref*{theory:ethics:organization}節}

\par 参考：\cite{fudano_ethics2021}。

\par 難易度 2（校準初稿）、優先度未評価。

\subsection{令和2年度 問2}
\label{q:r02-reg-apt-02}

\noindent\textbf{原典：} sources/R02/aptitude.pdf、PDF物理2頁。

\textbf{公式正答：} 2（single）。

\paragraph{問題の要旨} 学協会の倫理規程に沿う技術者の行動について、最も不適切な記述を選ぶ。

\Needspace{4\baselineskip}\paragraph{確認した条件・記述（独自要約）}\mbox{}\par

\noindent\textbf{方向：} 最も不適切な記述。\par

\paragraph{問われる核心} 事実に基づく報告と不正な例外の排除

\paragraph{必要知識} 真実性、公衆の優先、データの改ざん、組織の不正防止。

\paragraph{考え方・正答を選ぶ理由（査読待ち初稿）} 2は前半でデータの厳正な取扱いを述べながら、後半で顧客の要求に基づく不正な変更を許している。真実性は依頼者の都合によって解除されず、2を選ぶ。データを修正する行為全般が不正なのではなく、根拠のある訂正や処理は記録して説明する。

\Needspace{4\baselineskip}\paragraph{選択肢の検討（査読待ち初稿）}\mbox{}\par

\noindent\textbf{1：適切} 品質・信頼性・安全・環境への責任を持ち、公衆の安全・健康・福祉を最優先する。。専門職が負う技術的責任と公衆への優先的な配慮に沿う。\par

\noindent\textbf{2：不適切} 研究・調査の記録を厳正に扱い不正をしないが、顧客から要求があれば多少のデータ修正をしてよい。。顧客の望む結論に合わせた変更を正当化する例外はない。客観的な誤記訂正や妥当な処理とは区別する。\par

\noindent\textbf{3：適切} 属性による差別をせず、人権と人格を尊重する。。人格尊重と公平な取扱いに沿う。\par

\noindent\textbf{4：適切} 不正を防ぐ公正な環境の維持も責務とし、所属組織や技術者コミュニティの環境改善へ参加する。。不正防止を個人の心掛けだけに任せず、職務・研究環境を改善する責任を認めている。\par

\noindent\textbf{5：適切} 知識と経験を次世代の技術者・研究者の指導育成へ生かす。。人材育成と知識継承の責務に沿う。\par

\paragraph{再利用する考え方} 事実と価値判断、行為の結果と手続を区別し、事実に基づく報告と不正な例外の排除から判断の理由を説明する。

\paragraph{共通理論} \hyperref[theory:ethics:public]{第\ref*{theory:ethics:public}節}, \hyperref[theory:ethics:organization]{第\ref*{theory:ethics:organization}節}

\par 参考：\cite{ipej_ethics2011,mext_misconduct2014}。

\paragraph{関連問題} \hyperref[q:h24-reg-apt-02]{H24-REG-APT-02}

\par 難易度 2（校準初稿）、優先度未評価。

\subsection{令和5年度 問11}
\label{q:r05-reg-apt-11}

\noindent\textbf{原典：} sources/R05/aptitude.pdf、PDF物理11頁。

\textbf{公式正答：} 2（single）。

\paragraph{問題の要旨} エシックス・テストの4説明へ対応する名称の組合せを選ぶ。

\Needspace{4\baselineskip}\paragraph{確認した条件・記述（独自要約）}\mbox{}\par

\noindent\textbf{ア：} 全員が同じ行為をした場合の社会的成立と矛盾を考える。\par

\noindent\textbf{イ：} 影響を受ける側に自分を置いても同じ判断をするか考える。黄金律に基づく。\par

\noindent\textbf{ウ：} その行為を繰り返して選ぶと自分がどう見られるか考える。\par

\noindent\textbf{エ：} その行為が新聞等で報道された場合を考える。\par

\noindent\textbf{補足：} 専門家テストは専門職としてどう評価されるかを倫理綱領等で考える。\par

\noindent\textbf{並び：} ア・イ・ウ・エの順。\par

\paragraph{問われる核心} エシックス・テストの視点の違い

\paragraph{必要知識} 普遍化、黄金律と立場交換、人格、公開と説明可能性、専門職の規範。

\paragraph{考え方・正答を選ぶ理由（査読待ち初稿）} アは一般化したときの矛盾を調べる普遍化可能性、イは当事者の立場を逆にする可逆性である。ウは繰り返す行動が示す人格の美徳、エは公表されたとき説明できるかという世評で2を選ぶ。世評は単に人気を集めることではなく、判断を公開して吟味する道具である。

\Needspace{4\baselineskip}\paragraph{選択肢の検討（査読待ち初稿）}\mbox{}\par

\noindent\textbf{1：設問条件に不一致} 普遍化可能性／危害／世評／美徳。アだけは合うが、立場交換は可逆性、反復による人格は美徳、報道は世評である。\par

\noindent\textbf{2：設問条件に一致} 普遍化可能性／可逆性／美徳／世評。全員への適用、立場交換、反復による人格、公開時の評価がそれぞれ対応する。\par

\noindent\textbf{3：設問条件に不一致} 普遍化可能性／可逆性／世評／常識。アとイは合うが、ウは人格を問う美徳で、エは公開を問う世評である。\par

\noindent\textbf{4：設問条件に不一致} 常識／普遍化可能性／美徳／世評。ウとエは合うが、アは普遍化可能性、イは可逆性である。\par

\noindent\textbf{5：設問条件に不一致} 常識／危害／世評／普遍化可能性。全4対応が違う。危害の大小、一般化、立場交換、人格、公開は別の点検である。\par

\paragraph{再利用する考え方} 事実と価値判断、行為の結果と手続を区別し、エシックス・テストの視点の違いから判断の理由を説明する。

\paragraph{共通理論} \hyperref[theory:ethics:tests]{第\ref*{theory:ethics:tests}節}

\par 参考：\cite{davis_seven_steps}。

\par 難易度 2（校準初稿）、優先度未評価。

\subsection{令和7年度 問3}
\label{q:r07-reg-apt-03}

\noindent\textbf{原典：} sources/R07/aptitude.pdf、PDF物理3頁。

\textbf{公式正答：} 5（single）。

\paragraph{問題の要旨} 倫理的な行動・意思決定を阻害し得る5記述の数を選ぶ。

\Needspace{4\baselineskip}\paragraph{確認した条件・記述（独自要約）}\mbox{}\par

\noindent\textbf{導入例：} 私利私欲、責任を果たす結果への恐れ、他人のため・今回だけ等の自己欺瞞。\par

\noindent\textbf{ア：} 無知：適切な判断に不可欠な情報の知識がない。\par

\noindent\textbf{イ：} 自己中心的志向：自分の立場だけで判断し客観的な状況分析ができない。\par

\noindent\textbf{ウ：} 微視的視野：限定した範囲の詳細にとらわれる。\par

\noindent\textbf{エ：} 権威の無批判な受入れ：指導者や上司の意向・指示を吟味しない。\par

\noindent\textbf{オ：} 集団思考：十分な批判的考察なしに集団が結論に達する。\par

\paragraph{問われる核心} 個人と集団の判断を妨げる要因

\paragraph{必要知識} 必要情報、他者の視点、全体像、権威の点検、集団の批判的検討。

\paragraph{考え方・正答を選ぶ理由（査読待ち初稿）} アは情報不足、イは立場の偏り、ウは視野の狭さ、エは判断の放棄、オは集団での吟味不足を表す。いずれも阻害要因となり得るため5件で5を選ぶ。知識を得る努力、細部の調査、専門家への相談、協働そのものを否定する設問ではない。

\Needspace{4\baselineskip}\paragraph{選択肢の検討（査読待ち初稿）}\mbox{}\par

\noindent\textbf{1：設問条件に不一致} 1。不可欠な情報の欠如だけでなく、他の4項目も判断を妨げ得る。\par

\noindent\textbf{2：設問条件に不一致} 2。立場・視野・権威・集団の吟味不足も対象なので2件ではない。\par

\noindent\textbf{3：設問条件に不一致} 3。詳細への固執は細かく調べること一般と違い、周辺への影響を見落とすので3件ではない。\par

\noindent\textbf{4：設問条件に不一致} 4。集団で協力すること自体はよいが、批判的考察を省く集団思考も阻害要因である。\par

\noindent\textbf{5：設問条件に一致} 5。5項目全てが事実把握、他者への配慮、視野、自律、吟味を阻害し得る。\par

\paragraph{再利用する考え方} 事実と価値判断、行為の結果と手続を区別し、個人と集団の判断を妨げる要因から判断の理由を説明する。

\paragraph{共通理論} \hyperref[theory:ethics:obstacles]{第\ref*{theory:ethics:obstacles}節}

\par 参考：\cite{davis_seven_steps}。

\paragraph{関連問題} \hyperref[q:h29-reg-apt-15]{H29-REG-APT-15}

\par 難易度 2（校準初稿）、優先度未評価。

\section{倫理綱領の公共優先と誠実}

\subsection{平成28年度 問3}
\label{q:h28-reg-apt-03}

\noindent\textbf{原典：} sources/H28/aptitude.pdf、PDF物理3頁。

\textbf{公式正答：} 1（single）。

\paragraph{問題の要旨} 倫理綱領の公共優先・誠実・研鑽・知財。

\Needspace{4\baselineskip}\paragraph{確認した条件・記述（独自要約）}\mbox{}\par

\noindent\textbf{原典：} 問題文・否定・全五候補と組合せ表を保存原典36画像で確認。\par

\noindent\textbf{時点：} 平成28年度の本文表現を保持。後の制度改正を遡及しない。\par

\paragraph{問われる核心} 倫理綱領の公共優先・誠実・研鑽・知財。

\paragraph{必要知識} 倫理綱領の対象、客観性、説明責任、継続研鑽、知的財産。

\paragraph{考え方・正答を選ぶ理由（査読待ち初稿）} 公衆優先、事実・知識・良心による判断、継続研鑽、説明責任、他者の知的財産の尊重はすべて技術者倫理に適合する。OOOOOの1。専門職の信頼を能力だけで説明しない。

\Needspace{4\baselineskip}\paragraph{選択肢の検討（査読待ち初稿）}\mbox{}\par

\noindent\textbf{1：条件内で適切} OOOOO。各文又は定義の判定を先に行う。全項目が一致する。\par

\noindent\textbf{2：不適切} XOOOO。各文又は定義の判定を先に行う。並びが判定 OOOOO と異なる。\par

\noindent\textbf{3：不適切} OXOOO。各文又は定義の判定を先に行う。並びが判定 OOOOO と異なる。\par

\noindent\textbf{4：不適切} OOXOO。各文又は定義の判定を先に行う。並びが判定 OOOOO と異なる。\par

\noindent\textbf{5：不適切} OOOXO。各文又は定義の判定を先に行う。並びが判定 OOOOO と異なる。\par

\paragraph{再利用する考え方} 倫理綱領の対象、客観性、説明責任、継続研鑽、知的財産。 設問の否定・全称・基準年を確かめ、根拠と評価範囲を分けて判断する。

\paragraph{共通理論} \hyperref[theory:ethics:code]{第\ref*{theory:ethics:code}節}

\par 参考：\cite{ipej_ethics2011}。

\par 難易度 2（校準初稿）、優先度未評価。

\section{2023倫理綱領の指針と協働}

\subsection{令和6年度 問11}
\label{q:r06-reg-apt-11}

\noindent\textbf{原典：} sources/R06/aptitude.pdf、PDF物理11頁。

\textbf{公式正答：} 5（single）。

\paragraph{問題の要旨} 2023倫理綱領の五つの行動。

\Needspace{4\baselineskip}\paragraph{確認した条件・記述（独自要約）}\mbox{}\par

\noindent\textbf{原典：} 問題文・否定・全五候補と組合せ表を保存原典36画像で確認。\par

\noindent\textbf{時点：} 令和6年度の本文表現を保持。後の制度改正を遡及しない。\par

\paragraph{問われる核心} 2023倫理綱領の五つの行動。

\paragraph{必要知識} 2023綱領の指針、公共安全、持続可能性、秘密の対象、文化、人材育成。

\paragraph{考え方・正答を選ぶ理由（査読待ち初稿）} 5件すべて該当して5。リスクを根拠に検証し、環境・経済・社会への負の影響を減らし、正当な守秘と文化への尊重、人材育成を両立する。守秘を理由に公衆への責任まで消さない。

\Needspace{4\baselineskip}\paragraph{選択肢の検討（査読待ち初稿）}\mbox{}\par

\noindent\textbf{1：不適切} 1件。客観的リスク評価、負の影響の低減、正当な秘密の管理、現地文化の尊重、次世代育成はいずれも指針に含まれる。 判定を先に行った集計と一致しない。\par

\noindent\textbf{2：不適切} 2件。客観的リスク評価、負の影響の低減、正当な秘密の管理、現地文化の尊重、次世代育成はいずれも指針に含まれる。 判定を先に行った集計と一致しない。\par

\noindent\textbf{3：不適切} 3件。客観的リスク評価、負の影響の低減、正当な秘密の管理、現地文化の尊重、次世代育成はいずれも指針に含まれる。 判定を先に行った集計と一致しない。\par

\noindent\textbf{4：不適切} 4件。客観的リスク評価、負の影響の低減、正当な秘密の管理、現地文化の尊重、次世代育成はいずれも指針に含まれる。 判定を先に行った集計と一致しない。\par

\noindent\textbf{5：条件内で適切} 5件。客観的リスク評価、負の影響の低減、正当な秘密の管理、現地文化の尊重、次世代育成はいずれも指針に含まれる。 集計と一致する。\par

\paragraph{再利用する考え方} 2023綱領の指針、公共安全、持続可能性、秘密の対象、文化、人材育成。 設問の否定・全称・基準年を確かめ、根拠と評価範囲を分けて判断する。

\paragraph{共通理論} \hyperref[theory:ethics:code]{第\ref*{theory:ethics:code}節}

\par 参考：\cite{ipej_ethics2023}。

\paragraph{関連問題} \hyperref[q:r07-reg-apt-04]{R07-REG-APT-04}

\par 難易度 2（校準初稿）、優先度未評価。

\subsection{令和7年度 問4}
\label{q:r07-reg-apt-04}

\noindent\textbf{原典：} sources/R07/aptitude.pdf、PDF物理4頁。

\textbf{公式正答：} 3（single）。

\paragraph{問題の要旨} 協働の力を借りて確信の持てる業務を行う。

\Needspace{4\baselineskip}\paragraph{確認した条件・記述（独自要約）}\mbox{}\par

\noindent\textbf{原典：} 問題文・否定・全五候補と組合せ表を保存原典36画像で確認。\par

\noindent\textbf{時点：} 令和7年度の本文表現を保持。後の制度改正を遡及しない。\par

\paragraph{問われる核心} 協働の力を借りて確信の持てる業務を行う。

\paragraph{必要知識} 専門能力の限界、協働、代替案、持続可能性、文化と人材育成。

\paragraph{考え方・正答を選ぶ理由（査読待ち初稿）} 不適切3。自分一人の能力だけで業務を抱え込むことを有能性としない。必要に応じて協業者の力を借り、責任と役割を明確にして確信の持てる範囲で行う。

\Needspace{4\baselineskip}\paragraph{選択肢の検討（査読待ち初稿）}\mbox{}\par

\noindent\textbf{1：条件内で適切} 公衆へ危害がある場合の代替案。客観的評価と改善により適切な解決を図る。\par

\noindent\textbf{2：条件内で適切} 環境・経済・社会の負の影響の低減。持続可能な社会への責任に沿う。\par

\noindent\textbf{3：不適切} 協業が必要でも他者へ依頼しない。自分と協業者の力量を生かす考え方に反し、必要な専門性を得る機会を捨てる。\par

\noindent\textbf{4：条件内で適切} 業務地域の文化を尊重。国・地域の慣行や生活様式、宗教等を理解する。\par

\noindent\textbf{5：条件内で適切} 社会に貢献する技術者を育成。継続研鑽と人材育成の指針に沿う。\par

\paragraph{再利用する考え方} 専門能力の限界、協働、代替案、持続可能性、文化と人材育成。 設問の否定・全称・基準年を確かめ、根拠と評価範囲を分けて判断する。

\paragraph{共通理論} \hyperref[theory:ethics:code]{第\ref*{theory:ethics:code}節}

\par 参考：\cite{ipej_ethics2023}。

\paragraph{関連問題} \hyperref[q:r06-reg-apt-11]{R06-REG-APT-11}

\par 難易度 2（校準初稿）、優先度未評価。

\section{発言内容と利益関係の開示}

\subsection{平成28年度 問12}
\label{q:h28-reg-apt-12}

\noindent\textbf{原典：} sources/H28/aptitude.pdf、PDF物理12頁。

\textbf{公式正答：} 5（single）。

\paragraph{問題の要旨} 客観的な投稿と依頼者との関係の開示。

\Needspace{4\baselineskip}\paragraph{確認した条件・記述（独自要約）}\mbox{}\par

\noindent\textbf{原典：} 問題文・否定・全五候補と組合せ表を保存原典36画像で確認。\par

\noindent\textbf{時点：} 平成28年度の本文表現を保持。後の制度改正を遡及しない。\par

\paragraph{問われる核心} 客観的な投稿と依頼者との関係の開示。

\paragraph{必要知識} 利益相反、真実性、客観性、公共への説明、依頼者との関係。

\paragraph{考え方・正答を選ぶ理由（査読待ち初稿）} 最も適切5。廃棄物処理をめぐる誤解を正す記事でも、都合のよい事実だけを選ばず、B社のコンサルタントという関係を示す。発言の内容の真実性と発言者の立場は別々に確認する。

\Needspace{4\baselineskip}\paragraph{選択肢の検討（査読待ち初稿）}\mbox{}\par

\noindent\textbf{1：不適切} 関係者だから投稿を必ず控える。関係があること自体から発言禁止を導かない。\par

\noindent\textbf{2：不適切} 関係者だから必ず擁護記事を出す。依頼者への誠実さは一方的な宣伝の義務ではない。\par

\noindent\textbf{3：不適切} 学協会の有識者に必ず代筆を頼む。第三者の起用だけでは本人の利益関係の説明を代替しない。\par

\noindent\textbf{4：不適切} 会社の広報担当者に必ず代筆を頼む。広報への委任は客観性と関係開示の要件を消さない。\par

\noindent\textbf{5：条件内で適切} 真実を客観的に示し会社との関係を開示。読者が主張と利害関係を併せて評価できるようにする。\par

\paragraph{再利用する考え方} 利益相反、真実性、客観性、公共への説明、依頼者との関係。 設問の否定・全称・基準年を確かめ、根拠と評価範囲を分けて判断する。

\paragraph{共通理論} \hyperref[theory:ethics:disclosure]{第\ref*{theory:ethics:disclosure}節}

\par 参考：\cite{ipej_ethics2011,scu_ethics_framework}。

\par 難易度 2（校準初稿）、優先度未評価。

\section{DEIの多様性・公平性・包摂性}

\subsection{令和7年度 問6}
\label{q:r07-reg-apt-06}

\noindent\textbf{原典：} sources/R07/aptitude.pdf、PDF物理6頁。

\textbf{公式正答：} 5（single）。

\paragraph{問題の要旨} DEIの多様性・公平性・包摂性。

\Needspace{4\baselineskip}\paragraph{確認した条件・記述（独自要約）}\mbox{}\par

\noindent\textbf{原典：} 問題文・否定・全五候補と組合せ表を保存原典36画像で確認。\par

\noindent\textbf{時点：} 令和7年度の本文表現を保持。後の制度改正を遡及しない。\par

\paragraph{問われる核心} DEIの多様性・公平性・包摂性。

\paragraph{必要知識} Diversity、Equity、Inclusion、個別の条件への支援と能力発揮。

\paragraph{考え方・正答を選ぶ理由（査読待ち初稿）} アの宣言の目的、イの属性・知識・経験・専門性等の多様性、ウの一人ひとりの特性に応じた公平な環境・道具、エの違いを包摂して能力を発揮し新たな価値を作ること、はいずれも正で5。公平性は全員へ同じ道具を配ることだけではない。

\Needspace{4\baselineskip}\paragraph{選択肢の検討（査読待ち初稿）}\mbox{}\par

\noindent\textbf{1：不適切} XOOO。各文又は定義の判定を先に行う。並びが判定 OOOO と異なる。\par

\noindent\textbf{2：不適切} OXOO。各文又は定義の判定を先に行う。並びが判定 OOOO と異なる。\par

\noindent\textbf{3：不適切} OOXO。各文又は定義の判定を先に行う。並びが判定 OOOO と異なる。\par

\noindent\textbf{4：不適切} OOOX。各文又は定義の判定を先に行う。並びが判定 OOOO と異なる。\par

\noindent\textbf{5：条件内で適切} OOOO。各文又は定義の判定を先に行う。全項目が一致する。\par

\paragraph{再利用する考え方} Diversity、Equity、Inclusion、個別の条件への支援と能力発揮。 設問の否定・全称・基準年を確かめ、根拠と評価範囲を分けて判断する。

\paragraph{共通理論} \hyperref[theory:ethics:diversity]{第\ref*{theory:ethics:diversity}節}

\par 参考：\cite{eth54_dei2025}。

\par 難易度 2（校準初稿）、優先度未評価。

\section{構造の危険と情報開示}

\subsection{令和元年度再試験 問4}
\label{q:r01-ret-apt-04}

\noindent\textbf{原典：} sources/R01\_retest/aptitude.pdf、PDF物理4頁。

\textbf{公式正答：} 1（single）。

\par\textbf{検討未確定：} 倫理事例の判断であり、すべての秘密情報を即時公開する一般的な法的義務へ置き換えない。実際の緊急性・通報先は個別に評価する。

\paragraph{問題の要旨} 構造の危険を隠さず関係者の安全へつなぐ。

\Needspace{4\baselineskip}\paragraph{確認した条件・記述（独自要約）}\mbox{}\par

\noindent\textbf{状況：} 使用中の建物と周辺の何千人もの人々、構造計算の誤り、補強に3か月、暴風時の監視と避難が可能。\par

\noindent\textbf{争点：} 混乱を恐れて入居者へ知らせないことの妥当性。\par

\noindent\textbf{主体：} Aは構造設計者、Bは市の担当技術者。Bを市の名称としない。\par

\noindent\textbf{ア：} Aが守秘を理由に依頼人の不公表方針へ従う。\par

\noindent\textbf{イ：} Aが安全を優先して依頼人を説得。\par

\noindent\textbf{ウ：} Bが混乱を恐れて不公表方針へ従う。\par

\noindent\textbf{エ：} Bが公衆の安全を優先して依頼人を説得。\par

\noindent\textbf{オ：} Bが関係者の知る権利を重視して依頼人を説得。\par

\paragraph{問われる核心} 構造の危険を隠さず関係者の安全へつなぐ。

\paragraph{必要知識} 公衆の安全、誤りの是正、避難・監視計画、依頼者への働きかけ、必要な情報開示。

\paragraph{考え方・正答を選ぶ理由（査読待ち初稿）} アはAが守秘を理由に依頼人の不公表方針へ従うため不適切。Aは既に依頼人へ計算の見落としを報告している。イはAが公衆の安全を優先して依頼人を説得するため適切。ウは市の担当技術者Bが混乱を懸念して不公表に従うため不適切。エはBが公衆の安全を優先して、オはBが従業員・周辺住民の知る権利を重視して、いずれも依頼人を説得するため適切。XOXOOの1。実現可能な避難計画にも関係者が必要な情報を得ることが重要で、無制限の秘密公開の一般論にはしない。

\Needspace{4\baselineskip}\paragraph{選択肢の検討（査読待ち初稿）}\mbox{}\par

\noindent\textbf{1：条件内で適切} XOXOO。各文又は定義の判定を先に行う。全項目が一致する。\par

\noindent\textbf{2：不適切} OXOXO。各文又は定義の判定を先に行う。並びが判定 XOXOO と異なる。\par

\noindent\textbf{3：不適切} OOXOX。各文又は定義の判定を先に行う。並びが判定 XOXOO と異なる。\par

\noindent\textbf{4：不適切} XXOOO。各文又は定義の判定を先に行う。並びが判定 XOXOO と異なる。\par

\noindent\textbf{5：不適切} OOXXO。各文又は定義の判定を先に行う。並びが判定 XOXOO と異なる。\par

\paragraph{再利用する考え方} 公衆の安全、誤りの是正、避難・監視計画、依頼者への働きかけ、必要な情報開示。 設問の否定・全称・基準年を確かめ、根拠と評価範囲を分けて判断する。

\paragraph{共通理論} \hyperref[theory:ethics:disclosure]{第\ref*{theory:ethics:disclosure}節}

\par 参考：\cite{ipej_ethics2011,scu_ethics_framework}。

\par 難易度 3（校準初稿）、優先度未評価。

\section{多様性の属性・経験・働き方}

\subsection{令和3年度 問9}
\label{q:r03-reg-apt-09}

\noindent\textbf{原典：} sources/R03/aptitude.pdf、PDF物理9頁。

\textbf{公式正答：} 1（single）。

\par\textbf{検討未確定：} 2025宣言は概念の補足であり2021に公布済みだったとしない。十項目の当時の出典全体と定義境界は査読対象。

\paragraph{問題の要旨} 多様性を性別以外にも広げて考える。

\Needspace{4\baselineskip}\paragraph{確認した条件・記述（独自要約）}\mbox{}\par

\noindent\textbf{原典：} 問題文・否定・全五候補と組合せ表を保存原典36画像で確認。\par

\noindent\textbf{時点：} 令和3年度の本文表現を保持。後の制度改正を遡及しない。\par

\paragraph{問われる核心} 多様性を性別以外にも広げて考える。

\paragraph{必要知識} 属性と経験の多様性、公平な機会、参加を支える環境。

\paragraph{考え方・正答を選ぶ理由（査読待ち初稿）} 挙げられた10項目を性別や国籍だけへ限定せず、すべて多様性を考える視点として扱う。該当しない0件の1。異なる価値観を尊重することは、差別や危害を与える行動の無条件容認ではない。

\Needspace{4\baselineskip}\paragraph{選択肢の検討（査読待ち初稿）}\mbox{}\par

\noindent\textbf{1：条件内で適切} 0件。性別・年齢・人種・国籍・障害・性的指向・宗教等・価値観・経験・働き方の10項目はいずれも多様性の対象。 集計と一致する。\par

\noindent\textbf{2：不適切} 1件。性別・年齢・人種・国籍・障害・性的指向・宗教等・価値観・経験・働き方の10項目はいずれも多様性の対象。 判定を先に行った集計と一致しない。\par

\noindent\textbf{3：不適切} 2件。性別・年齢・人種・国籍・障害・性的指向・宗教等・価値観・経験・働き方の10項目はいずれも多様性の対象。 判定を先に行った集計と一致しない。\par

\noindent\textbf{4：不適切} 3件。性別・年齢・人種・国籍・障害・性的指向・宗教等・価値観・経験・働き方の10項目はいずれも多様性の対象。 判定を先に行った集計と一致しない。\par

\noindent\textbf{5：不適切} 4件。性別・年齢・人種・国籍・障害・性的指向・宗教等・価値観・経験・働き方の10項目はいずれも多様性の対象。 判定を先に行った集計と一致しない。\par

\paragraph{再利用する考え方} 属性と経験の多様性、公平な機会、参加を支える環境。 設問の否定・全称・基準年を確かめ、根拠と評価範囲を分けて判断する。

\paragraph{共通理論} \hyperref[theory:ethics:diversity]{第\ref*{theory:ethics:diversity}節}

\par 参考：\cite{eth54_dei2025}。

\par 難易度 3（校準初稿）、優先度未評価。

\section{ISO26000の七原則}

\subsection{平成29年度 問11}
\label{q:h29-reg-apt-11}

\noindent\textbf{原典：} sources/H29/aptitude.pdf、PDF物理11頁。

\textbf{公式正答：} 5（single）。

\paragraph{問題の要旨} ISO26000の七原則と技術の重要性を区別。

\Needspace{4\baselineskip}\paragraph{確認した条件・記述（独自要約）}\mbox{}\par

\noindent\textbf{原典：} 問題文・否定・全五候補と組合せ表を保存原典36画像で確認。\par

\noindent\textbf{時点：} 平成29年度の本文表現を保持。後の制度改正を遡及しない。\par

\paragraph{問われる核心} ISO26000の七原則と技術の重要性を区別。

\paragraph{必要知識} 説明責任、透明性、倫理的行動、ステークホルダーの利害の尊重、法の支配の尊重、国際行動規範の尊重、人権の尊重。

\paragraph{考え方・正答を選ぶ理由（査読待ち初稿）} 七原則に該当しない5。説明責任、透明性、倫理的行動、ステークホルダーの利害の尊重、法の支配の尊重、国際行動規範の尊重、人権の尊重。中核主題の七つと混同しない。重要な活動であっても規格の原則の見出しになるとは限らない。

\Needspace{4\baselineskip}\paragraph{選択肢の検討（査読待ち初稿）}\mbox{}\par

\noindent\textbf{1：条件内で適切} ステークホルダーの利害の尊重。七原則に含まれる。社会的影響を受ける者の利害を検討する。\par

\noindent\textbf{2：条件内で適切} 法の支配の尊重。七原則に含まれる。組織の独自ルールだけで法を置き換えない。\par

\noindent\textbf{3：条件内で適切} 国際行動規範の尊重。七原則に含まれる。国内の慣行だけへ限定しない。\par

\noindent\textbf{4：条件内で適切} 人権の尊重。七原則に含まれる。人権は中核主題にも関係する。\par

\noindent\textbf{5：不適切} 技術・ノウハウの重要性。技術が不要という意味ではなく、七原則の名称として挙げられていない。\par

\paragraph{再利用する考え方} 説明責任、透明性、倫理的行動、ステークホルダーの利害の尊重、法の支配の尊重、国際行動規範の尊重、人権の尊重。 設問の否定・全称・基準年を確かめ、根拠と評価範囲を分けて判断する。

\paragraph{共通理論} \hyperref[theory:ethics:iso]{第\ref*{theory:ethics:iso}節}

\par 参考：\cite{eth54_iso_brochure,iso26000}。

\paragraph{関連問題} \hyperref[q:r04-reg-apt-03]{R04-REG-APT-03}

\par 難易度 2（校準初稿）、優先度未評価。

\subsection{令和4年度 問3}
\label{q:r04-reg-apt-03}

\noindent\textbf{原典：} sources/R04/aptitude.pdf、PDF物理2頁。

\textbf{公式正答：} 3（single）。

\paragraph{問題の要旨} 社会的責任の七原則の名称。

\Needspace{4\baselineskip}\paragraph{確認した条件・記述（独自要約）}\mbox{}\par

\noindent\textbf{原典：} 問題文・否定・全五候補と組合せ表を保存原典36画像で確認。\par

\noindent\textbf{時点：} 令和4年度の本文表現を保持。後の制度改正を遡及しない。\par

\paragraph{問われる核心} 社会的責任の七原則の名称。

\paragraph{必要知識} 説明責任、透明性、倫理的行動、ステークホルダーの利害の尊重、法の支配の尊重、国際行動規範の尊重、人権の尊重。

\paragraph{考え方・正答を選ぶ理由（査読待ち初稿）} 七原則に該当しない3。説明責任、透明性、倫理的行動、ステークホルダーの利害の尊重、法の支配の尊重、国際行動規範の尊重、人権の尊重。技術継承の重要性を否定する問題ではない。

\Needspace{4\baselineskip}\paragraph{選択肢の検討（査読待ち初稿）}\mbox{}\par

\noindent\textbf{1：条件内で適切} 透明性。七原則に含まれる。影響の説明可能性を支える。\par

\noindent\textbf{2：条件内で適切} 法の支配の尊重。七原則に含まれる。適法性を基本として考える。\par

\noindent\textbf{3：不適切} 技術継承。重要な取組ではあるが七原則の見出しではない。\par

\noindent\textbf{4：条件内で適切} 説明責任。七原則に含まれる。影響と判断を説明する。\par

\noindent\textbf{5：条件内で適切} ステークホルダーの利害の尊重。七原則に含まれる。関係者の利害を考慮する。\par

\paragraph{再利用する考え方} 説明責任、透明性、倫理的行動、ステークホルダーの利害の尊重、法の支配の尊重、国際行動規範の尊重、人権の尊重。 設問の否定・全称・基準年を確かめ、根拠と評価範囲を分けて判断する。

\paragraph{共通理論} \hyperref[theory:ethics:iso]{第\ref*{theory:ethics:iso}節}

\par 参考：\cite{eth54_iso_brochure}。

\paragraph{関連問題} \hyperref[q:h29-reg-apt-11]{H29-REG-APT-11}, \hyperref[q:r06-reg-apt-04]{R06-REG-APT-04}

\par 難易度 2（校準初稿）、優先度未評価。

\subsection{令和6年度 問4}
\label{q:r06-reg-apt-04}

\noindent\textbf{原典：} sources/R06/aptitude.pdf、PDF物理4頁。

\textbf{公式正答：} 1（single）。

\par\textbf{検討未確定：} ISO公式要旨と七原則図を確認。JIS Z26000:2012の全翻訳本文の逐語照合は未完。

\paragraph{問題の要旨} 社会的責任の目的・規範・多様性・透明性・人権。

\Needspace{4\baselineskip}\paragraph{確認した条件・記述（独自要約）}\mbox{}\par

\noindent\textbf{原典：} 問題文・否定・全五候補と組合せ表を保存原典36画像で確認。\par

\noindent\textbf{時点：} 令和6年度の本文表現を保持。後の制度改正を遡及しない。\par

\paragraph{問われる核心} 社会的責任の目的・規範・多様性・透明性・人権。

\paragraph{必要知識} 社会的責任の目標、法令を基礎とする行動、国際規範と多様性、七原則。

\paragraph{考え方・正答を選ぶ理由（査読待ち初稿）} 目標は持続可能な発展、国際行動規範との整合を保って各種の多様性を考慮し、七原則の透明性と人権の尊重を読む。五語が揃う1。透明性は個人情報・正当な秘密まで無制限に公開する意味ではない。

\Needspace{4\baselineskip}\paragraph{選択肢の検討（査読待ち初稿）}\mbox{}\par

\noindent\textbf{1：条件内で適切} 持続可能な発展・国際行動規範・多様性・透明性・人権の尊重。各文又は定義の判定を先に行う。全項目が一致する。\par

\noindent\textbf{2：不適切} 公共の福祉・国内法・多様性・機密性・公共の福祉。各文又は定義の判定を先に行う。並びが判定 持続可能な発展・国際行動規範・多様性・透明性・人権の尊重 と異なる。\par

\noindent\textbf{3：不適切} 持続可能な発展・国際行動規範・独自性・透明性・人権の尊重。各文又は定義の判定を先に行う。並びが判定 持続可能な発展・国際行動規範・多様性・透明性・人権の尊重 と異なる。\par

\noindent\textbf{4：不適切} 公共の福祉・国内法・独自性・機密性・公共の福祉。各文又は定義の判定を先に行う。並びが判定 持続可能な発展・国際行動規範・多様性・透明性・人権の尊重 と異なる。\par

\noindent\textbf{5：不適切} 持続可能な発展・国際行動規範・多様性・機密性・人権の尊重。各文又は定義の判定を先に行う。並びが判定 持続可能な発展・国際行動規範・多様性・透明性・人権の尊重 と異なる。\par

\paragraph{再利用する考え方} 社会的責任の目標、法令を基礎とする行動、国際規範と多様性、七原則。 設問の否定・全称・基準年を確かめ、根拠と評価範囲を分けて判断する。

\paragraph{共通理論} \hyperref[theory:ethics:iso]{第\ref*{theory:ethics:iso}節}

\par 参考：\cite{iso26000,eth54_iso_brochure}。

\paragraph{関連問題} \hyperref[q:r04-reg-apt-03]{R04-REG-APT-03}

\par 難易度 3（校準初稿）、優先度未評価。

\section{規範の効力と専門職の責任}

\subsection{平成27年度 問5}
\label{q:h27-reg-apt-05}

\noindent\textbf{原典：} sources/H27/aptitude.pdf、PDF物理5頁。

\textbf{公式正答：} 5（single）。

\paragraph{問題の要旨} 倫理綱領・社会規範・専門職の責任。

\Needspace{4\baselineskip}\paragraph{確認した条件・記述（独自要約）}\mbox{}\par

\noindent\textbf{原典：} 問題文・否定・全五候補と組合せ表を保存原典36画像で確認。\par

\noindent\textbf{時点：} 平成27年度の本文表現を保持。後の制度改正を遡及しない。\par

\paragraph{問われる核心} 倫理綱領・社会規範・専門職の責任。

\paragraph{必要知識} 団体規範の効力、社会規範の種類、公衆の安全、法令遵守を含む倫理。

\paragraph{考え方・正答を選ぶ理由（査読待ち初稿）} 綱領の会員上の効力は所属団体に関係するが、非会員だから法や契約上の責任までなくなるとはしない。法律、契約、手順、綱領、慣習は規範の異なる形であり、専門職の公衆優先と広いコンプライアンスも適切。OOOOの5。

\Needspace{4\baselineskip}\paragraph{選択肢の検討（査読待ち初稿）}\mbox{}\par

\noindent\textbf{1：不適切} XOOO。各文又は定義の判定を先に行う。並びが判定 OOOO と異なる。\par

\noindent\textbf{2：不適切} OXOO。各文又は定義の判定を先に行う。並びが判定 OOOO と異なる。\par

\noindent\textbf{3：不適切} OOXO。各文又は定義の判定を先に行う。並びが判定 OOOO と異なる。\par

\noindent\textbf{4：不適切} OOOX。各文又は定義の判定を先に行う。並びが判定 OOOO と異なる。\par

\noindent\textbf{5：条件内で適切} OOOO。各文又は定義の判定を先に行う。全項目が一致する。\par

\paragraph{再利用する考え方} 団体規範の効力、社会規範の種類、公衆の安全、法令遵守を含む倫理。 設問の否定・全称・基準年を確かめ、根拠と評価範囲を分けて判断する。

\paragraph{共通理論} \hyperref[theory:ethics:code]{第\ref*{theory:ethics:code}節}

\par 参考：\cite{ipej_ethics2011}。

\par 難易度 2（校準初稿）、優先度未評価。

\section{功利・正義・義務の倫理視点}

\subsection{平成25年度 問15}
\label{q:h25-reg-apt-15}

\noindent\textbf{原典：} sources/H25/aptitude.pdf、PDF物理17頁。

\textbf{公式正答：} 3（single）。

\paragraph{問題の要旨} 功利・正義・義務の三つの倫理視点。

\Needspace{4\baselineskip}\paragraph{確認した条件・記述（独自要約）}\mbox{}\par

\noindent\textbf{原典：} 問題文・否定・全五候補と組合せ表を保存原典36画像で確認。\par

\noindent\textbf{時点：} 平成25年度の本文表現を保持。後の制度改正を遡及しない。\par

\paragraph{問われる核心} 功利・正義・義務の三つの倫理視点。

\paragraph{必要知識} 功利主義、分配の正義、義務と普遍化。

\paragraph{考え方・正答を選ぶ理由（査読待ち初稿）} 功利主義は結果の総利益を検討するb、正義の倫理は公平と弱い立場の権利を問うc、義務の倫理は普遍化できる行為規則を問うa。よってb c aの3。利益の合計だけで個人の権利が自動的に消えるとはしない。

\Needspace{4\baselineskip}\paragraph{選択肢の検討（査読待ち初稿）}\mbox{}\par

\noindent\textbf{1：不適切} a c b。各文又は定義の判定を先に行う。並びが判定 b c a と異なる。\par

\noindent\textbf{2：不適切} b a c。各文又は定義の判定を先に行う。並びが判定 b c a と異なる。\par

\noindent\textbf{3：条件内で適切} b c a。各文又は定義の判定を先に行う。全項目が一致する。\par

\noindent\textbf{4：不適切} c a b。各文又は定義の判定を先に行う。並びが判定 b c a と異なる。\par

\noindent\textbf{5：不適切} c b a。各文又は定義の判定を先に行う。並びが判定 b c a と異なる。\par

\paragraph{再利用する考え方} 功利主義、分配の正義、義務と普遍化。 設問の否定・全称・基準年を確かめ、根拠と評価範囲を分けて判断する。

\paragraph{共通理論} \hyperref[theory:ethics:principles]{第\ref*{theory:ethics:principles}節}

\par 参考：\cite{scu_ethics_framework}。

\paragraph{関連問題} \hyperref[q:r04-reg-apt-07]{R04-REG-APT-07}

\par 難易度 2（校準初稿）、優先度未評価。

\subsection{令和4年度 問7}
\label{q:r04-reg-apt-07}

\noindent\textbf{原典：} sources/R04/aptitude.pdf、PDF物理6頁。

\textbf{公式正答：} 3（single）。

\paragraph{問題の要旨} 功利・権利・黄金律・健康の判断語。

\Needspace{4\baselineskip}\paragraph{確認した条件・記述（独自要約）}\mbox{}\par

\noindent\textbf{原典：} 問題文・否定・全五候補と組合せ表を保存原典36画像で確認。\par

\noindent\textbf{時点：} 令和4年度の本文表現を保持。後の制度改正を遡及しない。\par

\paragraph{問われる核心} 功利・権利・黄金律・健康の判断語。

\paragraph{必要知識} 功利主義、個人の権利、相互性の黄金律、公衆の安全・健康・福利。

\paragraph{考え方・正答を選ぶ理由（査読待ち初稿）} 功利の判断は最大幸福、権利の判断は個人の権利、黄金律は自分が望む扱いを他者にも行う発想、公衆の優先対象は安全・健康・福利である。四語が揃う3。黄金律の適用でも相手の事情と価値観を確認する。

\Needspace{4\baselineskip}\paragraph{選択肢の検討（査読待ち初稿）}\mbox{}\par

\noindent\textbf{1：不適切} 最大幸福・多数派の権利・自分の望むことを人にせよ・身分。各文又は定義の判定を先に行う。並びが判定 最大幸福・個人の権利・自分の望むことを人にせよ・健康 と異なる。\par

\noindent\textbf{2：不適切} 最大利潤・個人の権利・人が望むことを自分にせよ・健康。各文又は定義の判定を先に行う。並びが判定 最大幸福・個人の権利・自分の望むことを人にせよ・健康 と異なる。\par

\noindent\textbf{3：条件内で適切} 最大幸福・個人の権利・自分の望むことを人にせよ・健康。各文又は定義の判定を先に行う。全項目が一致する。\par

\noindent\textbf{4：不適切} 最大利潤・多数派の権利・人が望むことを自分にせよ・健康。各文又は定義の判定を先に行う。並びが判定 最大幸福・個人の権利・自分の望むことを人にせよ・健康 と異なる。\par

\noindent\textbf{5：不適切} 最大幸福・個人の権利・人が望むことを自分にせよ・身分。各文又は定義の判定を先に行う。並びが判定 最大幸福・個人の権利・自分の望むことを人にせよ・健康 と異なる。\par

\paragraph{再利用する考え方} 功利主義、個人の権利、相互性の黄金律、公衆の安全・健康・福利。 設問の否定・全称・基準年を確かめ、根拠と評価範囲を分けて判断する。

\paragraph{共通理論} \hyperref[theory:ethics:principles]{第\ref*{theory:ethics:principles}節}

\par 参考：\cite{scu_ethics_framework,ipej_ethics2011}。

\paragraph{関連問題} \hyperref[q:h25-reg-apt-15]{H25-REG-APT-15}

\par 難易度 2（校準初稿）、優先度未評価。

\section{2007宣言の三つの責任}

\subsection{令和7年度 問2}
\label{q:r07-reg-apt-02}

\noindent\textbf{原典：} sources/R07/aptitude.pdf、PDF物理2頁。

\textbf{公式正答：} 1（single）。

\paragraph{問題の要旨} 2007プロフェッション宣言の三つの責任。

\Needspace{4\baselineskip}\paragraph{確認した条件・記述（独自要約）}\mbox{}\par

\noindent\textbf{原典：} 問題文・否定・全五候補と組合せ表を保存原典36画像で確認。\par

\noindent\textbf{時点：} 令和7年度の本文表現を保持。後の制度改正を遡及しない。\par

\paragraph{問われる核心} 2007プロフェッション宣言の三つの責任。

\paragraph{必要知識} 2007宣言、技術への責任、顧客と守秘義務、社会と公益。

\paragraph{考え方・正答を選ぶ理由（査読待ち初稿）} 自らの技術の向上に対する責任、守秘義務を守りつつ顧客へ誠実に履行する責任、社会・環境への影響と公益への責任を読む。四語が揃う1。顧客と社会の双方を考える点を残す。

\Needspace{4\baselineskip}\paragraph{選択肢の検討（査読待ち初稿）}\mbox{}\par

\noindent\textbf{1：条件内で適切} 技術・顧客・守秘義務・社会。各文又は定義の判定を先に行う。全項目が一致する。\par

\noindent\textbf{2：不適切} 技術・所属企業・基準・公衆。各文又は定義の判定を先に行う。並びが判定 技術・顧客・守秘義務・社会 と異なる。\par

\noindent\textbf{3：不適切} 行動・顧客・基準・社会。各文又は定義の判定を先に行う。並びが判定 技術・顧客・守秘義務・社会 と異なる。\par

\noindent\textbf{4：不適切} 行動・所属企業・守秘義務・公衆。各文又は定義の判定を先に行う。並びが判定 技術・顧客・守秘義務・社会 と異なる。\par

\noindent\textbf{5：不適切} 技術・顧客・基準・社会。各文又は定義の判定を先に行う。並びが判定 技術・顧客・守秘義務・社会 と異なる。\par

\paragraph{再利用する考え方} 2007宣言、技術への責任、顧客と守秘義務、社会と公益。 設問の否定・全称・基準年を確かめ、根拠と評価範囲を分けて判断する。

\paragraph{共通理論} \hyperref[theory:profession:declaration]{第\ref*{theory:profession:declaration}節}

\par 参考：\cite{eth54_profession2007}。

\par 難易度 2（校準初稿）、優先度未評価。

\section{公共責任と法の空白}

\subsection{平成25年度 問2}
\label{q:h25-reg-apt-02}

\noindent\textbf{原典：} sources/H25/aptitude.pdf、PDF物理3頁。

\textbf{公式正答：} 3（single）。

\paragraph{問題の要旨} 雇用者・顧客への誠実さと公共への責任。

\Needspace{4\baselineskip}\paragraph{確認した条件・記述（独自要約）}\mbox{}\par

\noindent\textbf{原典：} 問題文・否定・全五候補と組合せ表を保存原典36画像で確認。\par

\noindent\textbf{時点：} 平成25年度の本文表現を保持。後の制度改正を遡及しない。\par

\paragraph{問われる核心} 雇用者・顧客への誠実さと公共への責任。

\paragraph{必要知識} 専門職の複数の責任、公共安全、雇用者と顧客、研鑽と交流。

\paragraph{考え方・正答を選ぶ理由（査読待ち初稿）} アの法定責任を超える公共への責任、イの人材育成、エの継続研鑽と外部の専門的交流は適切。ウの「雇用者に対してのみ誠実」は顧客・公衆への責任を消すため誤。OOXOの3。

\Needspace{4\baselineskip}\paragraph{選択肢の検討（査読待ち初稿）}\mbox{}\par

\noindent\textbf{1：不適切} XOOO。各文又は定義の判定を先に行う。並びが判定 OOXO と異なる。\par

\noindent\textbf{2：不適切} OXOO。各文又は定義の判定を先に行う。並びが判定 OOXO と異なる。\par

\noindent\textbf{3：条件内で適切} OOXO。各文又は定義の判定を先に行う。全項目が一致する。\par

\noindent\textbf{4：不適切} OOOX。各文又は定義の判定を先に行う。並びが判定 OOXO と異なる。\par

\noindent\textbf{5：不適切} OOOO。各文又は定義の判定を先に行う。並びが判定 OOXO と異なる。\par

\paragraph{再利用する考え方} 専門職の複数の責任、公共安全、雇用者と顧客、研鑽と交流。 設問の否定・全称・基準年を確かめ、根拠と評価範囲を分けて判断する。

\paragraph{共通理論} \hyperref[theory:ethics:public]{第\ref*{theory:ethics:public}節}

\par 参考：\cite{ipej_ethics2011}。

\par 難易度 2（校準初稿）、優先度未評価。

\subsection{平成27年度 問4}
\label{q:h27-reg-apt-04}

\noindent\textbf{原典：} sources/H27/aptitude.pdf、PDF物理4頁。

\textbf{公式正答：} 3（single）。

\paragraph{問題の要旨} 法令の空白も公衆への配慮をなくさない。

\Needspace{4\baselineskip}\paragraph{確認した条件・記述（独自要約）}\mbox{}\par

\noindent\textbf{原典：} 問題文・否定・全五候補と組合せ表を保存原典36画像で確認。\par

\noindent\textbf{時点：} 平成27年度の本文表現を保持。後の制度改正を遡及しない。\par

\paragraph{問われる核心} 法令の空白も公衆への配慮をなくさない。

\paragraph{必要知識} 法の最低基準と倫理、事実の隠蔽、公共責任、適法な改善提案。

\paragraph{考え方・正答を選ぶ理由（査読待ち初稿）} アの法にないなら安全を考えなくてよい、イの不都合な知見を隠す、はいずれも誤。ウの法令の有無によらない公衆への責任と、エの合理性を欠く法令の改善を相談する行動は適切。XXOOの3。

\Needspace{4\baselineskip}\paragraph{選択肢の検討（査読待ち初稿）}\mbox{}\par

\noindent\textbf{1：不適切} OOXX。各文又は定義の判定を先に行う。並びが判定 XXOO と異なる。\par

\noindent\textbf{2：不適切} XOOO。各文又は定義の判定を先に行う。並びが判定 XXOO と異なる。\par

\noindent\textbf{3：条件内で適切} XXOO。各文又は定義の判定を先に行う。全項目が一致する。\par

\noindent\textbf{4：不適切} OOOX。各文又は定義の判定を先に行う。並びが判定 XXOO と異なる。\par

\noindent\textbf{5：不適切} XXOX。各文又は定義の判定を先に行う。並びが判定 XXOO と異なる。\par

\paragraph{再利用する考え方} 法の最低基準と倫理、事実の隠蔽、公共責任、適法な改善提案。 設問の否定・全称・基準年を確かめ、根拠と評価範囲を分けて判断する。

\paragraph{共通理論} \hyperref[theory:ethics:public]{第\ref*{theory:ethics:public}節}

\par 参考：\cite{ipej_ethics2011}。

\paragraph{関連問題} \hyperref[q:h28-reg-apt-09]{H28-REG-APT-09}

\par 難易度 2（校準初稿）、優先度未評価。

\subsection{平成28年度 問9}
\label{q:h28-reg-apt-09}

\noindent\textbf{原典：} sources/H28/aptitude.pdf、PDF物理9頁。

\textbf{公式正答：} 5（single）。

\paragraph{問題の要旨} 公共責任と不都合な事実の扱い。

\Needspace{4\baselineskip}\paragraph{確認した条件・記述（独自要約）}\mbox{}\par

\noindent\textbf{原典：} 問題文・否定・全五候補と組合せ表を保存原典36画像で確認。\par

\noindent\textbf{時点：} 平成28年度の本文表現を保持。後の制度改正を遡及しない。\par

\paragraph{問われる核心} 公共責任と不都合な事実の扱い。

\paragraph{必要知識} 法律と倫理、事実の報告、安全配慮、制度の改善。

\paragraph{考え方・正答を選ぶ理由（査読待ち初稿）} アの法定責任を超える公共配慮は正。イの法令の空白を理由に安全配慮を省くことは誤。ウの必要な法令の改善に向けた相談は正。エの不都合な知見の隠蔽は誤。OXOXの5。

\Needspace{4\baselineskip}\paragraph{選択肢の検討（査読待ち初稿）}\mbox{}\par

\noindent\textbf{1：不適切} OXXX。各文又は定義の判定を先に行う。並びが判定 OXOX と異なる。\par

\noindent\textbf{2：不適切} OXOO。各文又は定義の判定を先に行う。並びが判定 OXOX と異なる。\par

\noindent\textbf{3：不適切} XXOX。各文又は定義の判定を先に行う。並びが判定 OXOX と異なる。\par

\noindent\textbf{4：不適切} OOOX。各文又は定義の判定を先に行う。並びが判定 OXOX と異なる。\par

\noindent\textbf{5：条件内で適切} OXOX。各文又は定義の判定を先に行う。全項目が一致する。\par

\paragraph{再利用する考え方} 法律と倫理、事実の報告、安全配慮、制度の改善。 設問の否定・全称・基準年を確かめ、根拠と評価範囲を分けて判断する。

\paragraph{共通理論} \hyperref[theory:ethics:public]{第\ref*{theory:ethics:public}節}

\par 参考：\cite{ipej_ethics2011}。

\paragraph{関連問題} \hyperref[q:h27-reg-apt-04]{H27-REG-APT-04}

\par 難易度 2（校準初稿）、優先度未評価。

\section{公衆と情報に基づく同意}

\subsection{令和3年度 問2}
\label{q:r03-reg-apt-02}

\noindent\textbf{原典：} sources/R03/aptitude.pdf、PDF物理2頁。

\textbf{公式正答：} 4（single）。

\par\textbf{検討未確定：} 原典の公衆モデルをすべての法令や辞書の定義として扱わない。出典モデル原著の全章照合は未完。

\paragraph{問題の要旨} 倫理上の公衆は専門的な同意の条件から捉える。

\Needspace{4\baselineskip}\paragraph{確認した条件・記述（独自要約）}\mbox{}\par

\noindent\textbf{原典：} 問題文・否定・全五候補と組合せ表を保存原典36画像で確認。\par

\noindent\textbf{時点：} 令和3年度の本文表現を保持。後の制度改正を遡及しない。\par

\paragraph{問われる核心} 倫理上の公衆は専門的な同意の条件から捉える。

\paragraph{必要知識} 公衆、情報に基づく同意、専門知識の非対称性、説明責任。

\paragraph{考え方・正答を選ぶ理由（査読待ち初稿）} 設問の倫理的モデルでは、専門的な影響について十分な情報を得て同意することが難しい人々という4を選ぶ。能力や人格の優劣ではなく、技術知識と判断条件の非対称性を表す。

\Needspace{4\baselineskip}\paragraph{選択肢の検討（査読待ち初稿）}\mbox{}\par

\noindent\textbf{1：不適切} 国家や社会を形成する一般の人々。日常語としては使えるが、ここで問う専門知識と意思決定の条件を特定しない。\par

\noindent\textbf{2：不適切} 背景を異にする多数の組織されていない人々。集団の背景と組織化だけでは、影響と情報に基づく同意の条件を説明しない。\par

\noindent\textbf{3：不適切} 専門職としての技術業についていない人々。職業区分だけでは、その技術的影響について十分な情報を得て同意できるかを説明しない。\par

\noindent\textbf{4：条件内で適切} 十分な情報に基づく同意ができない人々。専門的な技術の影響に対し、知識や情報の不足から自律的判断が難しい立場を捉える。\par

\noindent\textbf{5：不適切} 広い地域に散在しメディアを通じ世論を形成する人々。世論形成の説明であり、技術的影響への informed consent の条件ではない。\par

\paragraph{再利用する考え方} 公衆、情報に基づく同意、専門知識の非対称性、説明責任。 設問の否定・全称・基準年を確かめ、根拠と評価範囲を分けて判断する。

\paragraph{共通理論} \hyperref[theory:ethics:public]{第\ref*{theory:ethics:public}節}

\par 参考：\cite{scu_ethics_framework,ipej_ethics2011}。

\par 難易度 3（校準初稿）、優先度未評価。

\section{組織の社会的責任と不正対応}

\subsection{平成26年度 問12}
\label{q:h26-reg-apt-12}

\noindent\textbf{原典：} sources/H26/aptitude.pdf、PDF物理12頁。

\textbf{公式正答：} 2（single）。

\par\textbf{検討未確定：} 大学の個別規程は一般的な反例であり、全大学・2014年の処分分布の統計を確認したものではない。

\paragraph{問題の要旨} 組織の社会的責任と大学の不正対応。

\Needspace{4\baselineskip}\paragraph{確認した条件・記述（独自要約）}\mbox{}\par

\noindent\textbf{原典：} 問題文・否定・全五候補と組合せ表を保存原典36画像で確認。\par

\noindent\textbf{時点：} 平成26年度の本文表現を保持。後の制度改正を遡及しない。\par

\paragraph{問われる核心} 組織の社会的責任と大学の不正対応。

\paragraph{必要知識} 大学の規程と不正対応、組織の説明責任、社会的責任の規格と自発的原則。

\paragraph{考え方・正答を選ぶ理由（査読待ち初稿）} 不適切2。大学の不正行為も教育・組織の公正を損ない、規程に応じた処分の対象となる。企業の責任、学問の自由、ISOの手引、国連の四分野は対象と性格を分けて理解する。

\Needspace{4\baselineskip}\paragraph{選択肢の検討（査読待ち初稿）}\mbox{}\par

\noindent\textbf{1：条件内で適切} 企業の社会的責任を積極的に果たす。規則の最低限だけでなく社会との信頼関係を含む。\par

\noindent\textbf{2：不適切} 大学の試験不正は注意だけで単位処分なしが一般。大学でも学則等に基づく懲戒や科目の成績措置がある。注意のみを一般化できない。\par

\noindent\textbf{3：条件内で適切} 学問の自由と組織の健全性を両立。自由は不正な評価や責任の不存在を意味しない。\par

\noindent\textbf{4：条件内で適切} ISO26000は持続可能な発展へ貢献。認証要求ではなく社会的責任の手引である。\par

\noindent\textbf{5：条件内で適切} 国連グローバル・コンパクトの四分野。人権、労働、環境、腐敗防止という対象分野を区別する。\par

\paragraph{再利用する考え方} 大学の規程と不正対応、組織の説明責任、社会的責任の規格と自発的原則。 設問の否定・全称・基準年を確かめ、根拠と評価範囲を分けて判断する。

\paragraph{共通理論} \hyperref[theory:ethics:csr]{第\ref*{theory:ethics:csr}節}

\par 参考：\cite{iso26000,eth54_globalcompact,eth54_university}。

\par 難易度 3（校準初稿）、優先度未評価。

\section{非営利活動の公益目的・事業収入と報酬}

\subsection{平成24年度 問13}
\label{q:h24-reg-apt-13}

\noindent\textbf{原典：} sources/H24/aptitude.pdf、PDF物理11頁。

\textbf{公式正答：} 3（single）。

\paragraph{問題の要旨} 特定非営利活動法人の公益目的、収益・報酬、競争、会社員の参加を4記述で判定する。

\Needspace{4\baselineskip}\paragraph{確認した条件・記述（独自要約）}\mbox{}\par

\noindent\textbf{方向：} ア〜エの正誤組合せ。\par

\noindent\textbf{時点：} 平成24年、2012年4月施行の改正後。\par

\noindent\textbf{限定：} イは法人の目的と利益分配の制限の文脈。無報酬や全ての個人への便益禁止へ拡張しない。\par

\paragraph{問われる核心} 非営利の公益目的と事業収入・労働報酬の区別

\paragraph{必要知識} 非営利の公益目的と事業収入・労働報酬の区別。アは○。2012年4月施行の改正後、法定の20分野と不特定多数の利益という要件を照合する。イは○。非営利とは構成員への利益分配を目的としないことで、収入・剰余・給与が全くない意味ではない。ただし役員として報酬を受ける者は役員総数の3分の1以下などの条件がある。ウは×。法定分野の制約はあるが、営利企業との競争を保護するために活動範囲を制限するという根拠ではない。エは×。就業規則等で認める会社員の参加をNPOという形態自体が一律に禁ずることはない。○○××の3。

\paragraph{考え方・正答を選ぶ理由（査読待ち初稿）} アは○。2012年4月施行の改正後、法定の20分野と不特定多数の利益という要件を照合する。イは○。非営利とは構成員への利益分配を目的としないことで、収入・剰余・給与が全くない意味ではない。ただし役員として報酬を受ける者は役員総数の3分の1以下などの条件がある。ウは×。法定分野の制約はあるが、営利企業との競争を保護するために活動範囲を制限するという根拠ではない。エは×。就業規則等で認める会社員の参加をNPOという形態自体が一律に禁ずることはない。○○××の3。

\Needspace{4\baselineskip}\paragraph{選択肢の検討（査読待ち初稿）}\mbox{}\par

\noindent\textbf{1：選択対象ではない} ア○・イ○・ウ○・エ×。ウを適切とするのが誤り。活動分野の法定と営利企業の競争保護のための制限は別である。\par

\noindent\textbf{2：選択対象ではない} ア×・イ○・ウ○・エ×。アを誤り、ウを正しいとした点が逆。2012年4月以降は20分野である。\par

\noindent\textbf{3：選択対象} ア○・イ○・ウ×・エ×。公益目的と報酬の区別に基づくア・イの肯定、ウ・エの否定が一致する。\par

\noindent\textbf{4：選択対象ではない} ア○・イ○・ウ×・エ○。エを適切とするのが誤り。営利企業の社員であること自体は参加を禁ずる理由にならない。\par

\noindent\textbf{5：選択対象ではない} ア○・イ×・ウ×・エ×。イを否定した点が誤り。支障のないその他の事業と、役職員への適法な報酬は可能。\par

\paragraph{再利用する考え方} 対象、条件、時点を先に確認する。非営利の公益目的と事業収入・労働報酬の区別を具体的な事例へ適用し、技術的原因・法的要件・倫理的判断のどれを問うか区別する。

\paragraph{共通理論} \hyperref[theory:ethics:npo]{第\ref*{theory:ethics:npo}節}, \hyperref[theory:ethics:csr]{第\ref*{theory:ethics:csr}節}

\par 参考：\cite{cabinet_npo_overview,cabinet_npo_officers}。

\par 難易度 2（校準初稿）、優先度未評価。

\section{倫理的告発の証拠・段階・改善見込み}

\subsection{令和元年度再試験 問5}
\label{q:r01-ret-apt-05}

\noindent\textbf{原典：} sources/R01\_retest/aptitude.pdf、PDF物理5頁。

\textbf{公式正答：} 3（single）。

\par\textbf{検討未確定：} 倫理モデルの条件は実際の全緊急事態の手順を固定する法律ではない。直接の原著全章は未通読。急迫危険では別の法定経路と対応を検討する。

\paragraph{問題の要旨} 倫理的な警笛鳴らしの段階・証拠・改善見込み。

\Needspace{4\baselineskip}\paragraph{確認した条件・記述（独自要約）}\mbox{}\par

\noindent\textbf{原典：} 問題文・否定・全五候補と組合せ表を保存原典27画像で確認。\par

\noindent\textbf{時点：} 令和元年度再試験の本文表現を保持。後の制度改正を遡及しない。\par

\paragraph{問われる核心} 倫理的な警笛鳴らしの段階・証拠・改善見込み。

\paragraph{必要知識} 深刻な危害、内部の段階、客観的証拠、改善効果、個人の危険と公益。

\paragraph{考え方・正答を選ぶ理由（査読待ち初稿）} アは深刻な危害を認識した段階の内部報告で正。イは上司が反応しないだけで他の内部経路を検討せず即座に外へという飛躍、ウは客観的証拠を持たず外部告発義務を断定して誤。エは公表で必要な改善が見込める理由を求め正、3。これは倫理枠組みの比較で、法定保護の要件と一つにしない。

\Needspace{4\baselineskip}\paragraph{選択肢の検討（査読待ち初稿）}\mbox{}\par

\noindent\textbf{1：不適切} XOXO。各文又は定義の判定を先に行う。並びが判定 OXXO と異なる。\par

\noindent\textbf{2：不適切} OXOX。各文又は定義の判定を先に行う。並びが判定 OXXO と異なる。\par

\noindent\textbf{3：条件内で適切} OXXO。各文又は定義の判定を先に行う。全項目が一致する。\par

\noindent\textbf{4：不適切} XXOO。各文又は定義の判定を先に行う。並びが判定 OXXO と異なる。\par

\noindent\textbf{5：不適切} OOXX。各文又は定義の判定を先に行う。並びが判定 OXXO と異なる。\par

\paragraph{再利用する考え方} 対象・責任主体・時点・要件を分け、文の否定と「全て」「不要」の範囲を点検する。知識を当てはめる前に何を証明する資料かを確かめる。

\paragraph{共通理論} \hyperref[theory:ethics:whistle_conditions]{第\ref*{theory:ethics:whistle_conditions}節}

\par 参考：\cite{muroran_ethics_cases}。

\par 難易度 3（校準初稿）、優先度未評価。

\section{社会的責任・安全の標準}

\subsection{平成23年度 問11}
\label{q:h23-reg-apt-11}

\noindent\textbf{原典：} sources/H23/aptitude.pdf、PDF物理12頁。

\textbf{公式正答：} 3（single）。

\par\textbf{検討未確定：} Astra Review 1後の校正・解説例。解説はAstra査読待ち。①政府除外、④原則列挙、②⑤歴史的経緯を要再確認。要旨のみ校正済みで全文の校正完了ではない。

\paragraph{問題の要旨} 社会的責任の手引きと認証規格の区別を問う。公式正答③を保持し、他の記述の解釈上の疑義を別に残す。

\Needspace{4\baselineskip}\paragraph{確認した条件・記述（独自要約）}\mbox{}\par

\noindent\textbf{standard：} ISO26000\par

\noindent\textbf{exam\_reference\_year：} 2011\par

\paragraph{問われる核心} 技術者倫理・意思決定

\paragraph{必要知識} 社会的責任、要求事項とガイダンス、認証の適用範囲

\paragraph{考え方（査読待ち初稿）} ISO26000は社会的責任に関するガイダンスであり、認証のための要求事項を規定した規格ではない。③がこの性格を誤っているため公式正答は③。ISOの一次資料でも認証用途ではないことを確認した。

\paragraph{他の選択肢（査読待ち初稿）} 他肢をすべて無条件に正しいと説明しない。①の政府除外の表現、④の原則の列挙範囲には一次資料との照合を要する疑義が残る。②⑤の歴史的な参加・連携の詳細も当時資料の確認を後工程に残す。

\par 参考：\cite{iso26000}。

\par 難易度 3（校準初稿）、優先度未評価。

\subsection{平成24年度 問12}
\label{q:h24-reg-apt-12}

\noindent\textbf{原典：} sources/H24/aptitude.pdf、PDF物理10頁。

\textbf{公式正答：} 2（single）。

\paragraph{問題の要旨} 企業の社会的責任に関わる活動の捉え方について、最も不適切なものを選ぶ。

\Needspace{4\baselineskip}\paragraph{確認した条件・記述（独自要約）}\mbox{}\par

\noindent\textbf{用語：} CSRは企業の社会的責任、CSVは共有価値の創出。\par

\noindent\textbf{時点：} 平成24年度の設問。CSVは2011年の提案を背景とする。\par

\paragraph{問われる核心} 社会的価値と経済的価値の両立

\paragraph{必要知識} 社会的責任、利害関係者、寄付・事業活動・共有価値の区別。

\paragraph{考え方・正答を選ぶ理由（査読待ち初稿）} 企業の利益が生じるかどうかだけで社会的責任の活動を判別できない。2の「売上が上がるとCSRではなくなる」という排他的な条件が誤りである。5は社会と企業の両方の価値を生む戦略の提唱について述べるもので、全てのCSRがCSVになるとの意味ではない。

\Needspace{4\baselineskip}\paragraph{選択肢の検討（査読待ち初稿）}\mbox{}\par

\noindent\textbf{1：適切} 災害時に社員へ交通費を支給してボランティアとして派遣するのはCSRの一環となり得る。。社員の参加を支援して被災社会へ貢献する活動は、社会的責任の一つの取組となる。\par

\noindent\textbf{2：不適切} 社会貢献の結果売上が上がれば販売促進と見なされ、CSRにはなり得ない。。社会的効果と企業の経済的効果は両立し得る。売上増加だけを理由にCSRから排除できない。\par

\noindent\textbf{3：適切} 国内販売量に応じ海外経済支援へ資金提供するのはCSRの一環となり得る。。受益地域が販売地域と一致しなければならないという条件はない。\par

\noindent\textbf{4：適切} 自社製品の廃棄物を有償回収・再資源化して削減するのはCSRの一環となり得る。。料金を取ること自体は廃棄物削減の社会的効果を否定しない。\par

\noindent\textbf{5：適切} 受動的な対応にとどまらず企業戦略として共有価値の創出へ取り組む考え方が提唱される。。社会課題への対応と競争力を結び付けるCSVの提案は出題時点に存在する。ただしCSVとCSRを同一の定義にしない。\par

\paragraph{再利用する考え方} 事実と価値判断、行為の結果と手続を区別し、社会的価値と経済的価値の両立から判断の理由を説明する。

\paragraph{共通理論} \hyperref[theory:ethics:csr]{第\ref*{theory:ethics:csr}節}

\par 参考：\cite{iso26000,hbs_csv}。

\par 難易度 2（校準初稿）、優先度未評価。

\subsection{平成25年度 問6}
\label{q:h25-reg-apt-06}

\noindent\textbf{原典：} sources/H25/aptitude.pdf、PDF物理7頁。

\textbf{公式正答：} 5（single）。

\paragraph{問題の要旨} ISO26000の適用対象、持続可能な発展、利害関係者、認証用途を4記述で判定する。

\Needspace{4\baselineskip}\paragraph{確認した条件・記述（独自要約）}\mbox{}\par

\noindent\textbf{方向：} ア〜エの正誤組合せ。\par

\noindent\textbf{時点：} 2010年発行のISO26000、平成25年度。\par

\noindent\textbf{対象：} 組織の社会的責任。個人の責任一般を否定する設問ではない。\par

\paragraph{問われる核心} 社会的責任の対象・利害関係者と認証の区別

\paragraph{必要知識} 社会的責任の対象・利害関係者と認証の区別。ア×：ISO26000の適用対象は種類・規模を問わない組織であり、公人個人を別の独立した対象として加える説明は適切でない。個人の行為が組織の責任に関係することとは区別する。イ○：持続可能な発展への貢献が包括的目的である。ウ×：組織の決定・活動に利害を持つ従業員、消費者、地域等も関係し、株主だけに限定しない。エ×：要求事項を示す認証用規格ではなく、指針を示す。×○××の5。

\paragraph{考え方・正答を選ぶ理由（査読待ち初稿）} ア×：ISO26000の適用対象は種類・規模を問わない組織であり、公人個人を別の独立した対象として加える説明は適切でない。個人の行為が組織の責任に関係することとは区別する。イ○：持続可能な発展への貢献が包括的目的である。ウ×：組織の決定・活動に利害を持つ従業員、消費者、地域等も関係し、株主だけに限定しない。エ×：要求事項を示す認証用規格ではなく、指針を示す。×○××の5。

\Needspace{4\baselineskip}\paragraph{選択肢の検討（査読待ち初稿）}\mbox{}\par

\noindent\textbf{1：選択対象ではない} ア○・イ○・ウ×・エ×。イ・ウ・エは合うが、アは組織と公人個人を並列の適用対象とする点が誤り。\par

\noindent\textbf{2：選択対象ではない} ア○・イ×・ウ×・エ○。ア・イ・エの判断が誤り。発展への貢献は目的で、認証用ではない。\par

\noindent\textbf{3：選択対象ではない} ア×・イ×・ウ○・エ○。ア以外が誤り。株主のみへの限定、外部認証用という位置付けはいずれも違う。\par

\noindent\textbf{4：選択対象ではない} ア×・イ○・ウ×・エ○。ア・イ・ウは合うが、エの認証用という肯定が誤り。\par

\noindent\textbf{5：選択対象} ア×・イ○・ウ×・エ×。対象は組織、目的は持続可能な発展への貢献、株主限定でも認証用でもないという判断が全て一致。\par

\paragraph{再利用する考え方} 対象、条件、時点を先に確認する。社会的責任の対象・利害関係者と認証の区別を具体的な事例へ適用し、技術的原因・法的要件・倫理的判断のどれを問うか区別する。

\paragraph{共通理論} \hyperref[theory:ethics:iso]{第\ref*{theory:ethics:iso}節}, \hyperref[theory:ethics:csr]{第\ref*{theory:ethics:csr}節}

\par 参考：\cite{iso26000,iso26000_project2010}。

\paragraph{関連問題} \hyperref[q:r01-reg-apt-14]{R01-REG-APT-14}

\par 難易度 2（校準初稿）、優先度未評価。

\subsection{令和元年度 問14}
\label{q:r01-reg-apt-14}

\noindent\textbf{原典：} sources/R01/aptitude.pdf、PDF物理14頁。

\textbf{公式正答：} 4（single）。

\paragraph{問題の要旨} ISO26000の社会的責任の原則として、説明責任、透明性、倫理、法の支配、人権を判定する。

\Needspace{4\baselineskip}\paragraph{確認した条件・記述（独自要約）}\mbox{}\par

\noindent\textbf{方向：} 最も不適切。\par

\noindent\textbf{時点：} ISO26000:2010、JIS Z26000:2012を背景とする令和元年。\par

\paragraph{問われる核心} 社会的責任における法の支配と政府の意向の区別

\paragraph{必要知識} 社会的責任における法の支配と政府の意向の区別。七つの原則は、説明責任、透明性、倫理的行動、利害関係者の利害の尊重、法の支配の尊重、国際行動規範の尊重、人権の尊重である。4の後半はこの原則を政府の意向へ置き換える点が不適切である。七つの中核主題（組織統治、人権、労働慣行等）と原則のリストも別に区別する。

\paragraph{考え方・正答を選ぶ理由（査読待ち初稿）} 七つの原則は、説明責任、透明性、倫理的行動、利害関係者の利害の尊重、法の支配の尊重、国際行動規範の尊重、人権の尊重である。4の後半はこの原則を政府の意向へ置き換える点が不適切である。七つの中核主題（組織統治、人権、労働慣行等）と原則のリストも別に区別する。

\Needspace{4\baselineskip}\paragraph{選択肢の検討（査読待ち初稿）}\mbox{}\par

\noindent\textbf{1：適切} 組織は社会・経済・環境への影響に説明責任を負う。。行動の影響に応答し説明する原則に対応する。\par

\noindent\textbf{2：適切} 社会・環境に影響する決定と活動を透明にする。。判断と活動を説明可能・理解可能にする原則である。無関係な秘密まで無条件に公開する意味ではない。\par

\noindent\textbf{3：適切} 組織は倫理的に行動する。。誠実さ、公平さ等を行動に反映する原則に対応する。\par

\noindent\textbf{4：不適切} 法の支配を尊重すると同時に自国政府の意向も尊重する。。法の支配の尊重を政府の意向への同調へ拡張している。政府も法の制約を受け、意向は法・人権に優先する一般的根拠ではない。\par

\noindent\textbf{5：適切} 人権の重要性と普遍性を認識し尊重する。。対象国の都合だけで人権の普遍性を否定しない。\par

\paragraph{再利用する考え方} 対象、条件、時点を先に確認する。社会的責任における法の支配と政府の意向の区別を具体的な事例へ適用し、技術的原因・法的要件・倫理的判断のどれを問うか区別する。

\paragraph{共通理論} \hyperref[theory:ethics:iso]{第\ref*{theory:ethics:iso}節}

\par 参考：\cite{iso26000,iso_gri2014}。

\paragraph{関連問題} \hyperref[q:h25-reg-apt-06]{H25-REG-APT-06}

\par 難易度 2（校準初稿）、優先度未評価。

\section{ユニバーサルデザイン}

\subsection{令和2年度 問11}
\label{q:r02-reg-apt-11}

\noindent\textbf{原典：} sources/R02/aptitude.pdf、PDF物理11頁。

\textbf{公式正答：} 1（single）。

\paragraph{問題の要旨} 年齢・能力・言語等の多様性に配慮するユニバーサルデザインの7原則から、趣旨の違う記述の数を求める。

\Needspace{4\baselineskip}\paragraph{確認した条件・記述（独自要約）}\mbox{}\par

\noindent\textbf{方向：} 七原則と趣旨が異なる記述の数。候補0/1/2/3/4。\par

\noindent\textbf{7記述：} ア：公平な利用；イ：利用方法の柔軟性；ウ：単純で直観的な使用；エ：必要な情報の認知；オ：エラーの害を抑える設計；カ：少ない身体的努力；キ：接近・使用の寸法と空間\par

\noindent\textbf{時点：} 原典の国際イベント開催予定は令和2年出題当時の背景。現在の予定へ流用しない。\par

\paragraph{問われる核心} 利用者の多様性を設計段階から扱う社会的包摂

\paragraph{必要知識} 利用者の多様性を設計段階から扱う社会的包摂。ア〜キは順に、公平な利用、利用方法の柔軟性、単純で直観的な使用、必要な情報の認知、エラーの害を抑える設計、少ない身体的努力、接近・使用の寸法と空間であり、ノースカロライナ州立大学の1997年版の七原則へ全て対応する。したがって異なる趣旨は0項目で1を選ぶ。オのエラーへの寛容はミスを奨励する意味ではなく、偶発的な操作が重大な危害へ直結することを抑える。既存の障壁を除く対策も必要であり、設問前文からバリアフリーが常に過剰投資であるとの一般則を導かない。

\paragraph{考え方・正答を選ぶ理由（査読待ち初稿）} ア〜キは順に、公平な利用、利用方法の柔軟性、単純で直観的な使用、必要な情報の認知、エラーの害を抑える設計、少ない身体的努力、接近・使用の寸法と空間であり、ノースカロライナ州立大学の1997年版の七原則へ全て対応する。したがって異なる趣旨は0項目で1を選ぶ。オのエラーへの寛容はミスを奨励する意味ではなく、偶発的な操作が重大な危害へ直結することを抑える。既存の障壁を除く対策も必要であり、設問前文からバリアフリーが常に過剰投資であるとの一般則を導かない。

\Needspace{4\baselineskip}\paragraph{選択肢の検討（査読待ち初稿）}\mbox{}\par

\noindent\textbf{1：選択対象} 0項目。列挙された7項目は全て原則の趣旨に対応し、異なるものは0。\par

\noindent\textbf{2：選択対象ではない} 1項目。一つの原則を除外する根拠がない。\par

\noindent\textbf{3：選択対象ではない} 2項目。二つの原則を除外する根拠がない。\par

\noindent\textbf{4：選択対象ではない} 3項目。三つの原則を除外する根拠がない。\par

\noindent\textbf{5：選択対象ではない} 4項目。四つの原則を除外する根拠がない。\par

\paragraph{再利用する考え方} 対象、条件、時点を先に確認する。利用者の多様性を設計段階から扱う社会的包摂を具体的な事例へ適用し、技術的原因・法的要件・倫理的判断のどれを問うか区別する。

\paragraph{共通理論} \hyperref[theory:ethics:universal]{第\ref*{theory:ethics:universal}節}

\par 参考：\cite{ncsu_ud1997}。

\par 難易度 2（校準初稿）、優先度未評価。

\section{仕事と生活の調和・多様な働き方}

\subsection{平成27年度 問15}
\label{q:h27-reg-apt-15}

\noindent\textbf{原典：} sources/H27/aptitude.pdf、PDF物理16頁。

\textbf{公式正答：} 4（single）。

\paragraph{問題の要旨} 仕事と生活の調和の効果、継続改善、雇用形態、社会全体の協働を判定する。

\Needspace{4\baselineskip}\paragraph{確認した条件・記述（独自要約）}\mbox{}\par

\noindent\textbf{方向：} 最も不適切。\par

\noindent\textbf{時点：} 2007年策定、2010年改定の憲章を背景とする平成27年度。\par

\noindent\textbf{限定：} 効果は期待であり、個別企業への保証ではない。\par

\paragraph{問われる核心} 多様な選択を支える仕事と生活の調和

\paragraph{必要知識} 多様な選択を支える仕事と生活の調和。仕事の責任を果たしつつ人生の段階に応じた生活を実現することが目的である。短時間正社員や柔軟な勤務等も選択肢になり、正社員を減らすことへ限定する根拠はない。4が不適切。1・2は期待効果、3・5は改善の進め方として適切である。

\paragraph{考え方・正答を選ぶ理由（査読待ち初稿）} 仕事の責任を果たしつつ人生の段階に応じた生活を実現することが目的である。短時間正社員や柔軟な勤務等も選択肢になり、正社員を減らすことへ限定する根拠はない。4が不適切。1・2は期待効果、3・5は改善の進め方として適切である。

\Needspace{4\baselineskip}\paragraph{選択肢の検討（査読待ち初稿）}\mbox{}\par

\noindent\textbf{1：適切} 育児・介護・地域での役割を担いながら働け、労働力不足の緩和が期待できる。。働き続けることを妨げる条件を緩和する期待効果である。単一施策で人口問題が必ず解決するという断定ではない。\par

\noindent\textbf{2：適切} 企業の社会的責任の実践による信用・人材の確保や定着の改善が期待できる。。継続就業の環境は採用・定着にも関係する。全企業で直ちに同じ成果が出るとする保証ではない。\par

\noindent\textbf{3：適切} 多様な価値観・働き方へ、社会環境を継続的に改善する必要がある。。選好やライフステージが異なるため、単一の働き方の押し付けでは満たせない。制度と運用を見直す。\par

\noindent\textbf{4：不適切} 多様な働き方のために、正社員を減らし短期契約の非正規社員を増やす必要がある。。柔軟な時間・場所・休暇・職務設計と、雇用の不安定化は別の軸である。短期非正規への置換は必要条件ではなく、安定した生活を損なうこともある。\par

\noindent\textbf{5：適切} 経済界・労働界・国・自治体等が協力する必要がある。。職場だけでなく社会基盤や家庭事情にも依存するため、関係主体の役割分担が必要。\par

\paragraph{再利用する考え方} 対象、条件、時点を先に確認する。多様な選択を支える仕事と生活の調和を具体的な事例へ適用し、技術的原因・法的要件・倫理的判断のどれを問うか区別する。

\paragraph{共通理論} \hyperref[theory:ethics:worklife]{第\ref*{theory:ethics:worklife}節}, \hyperref[theory:ethics:organization]{第\ref*{theory:ethics:organization}節}

\par 参考：\cite{cabinet_worklife}。

\par 難易度 2（校準初稿）、優先度未評価。
